% Chapter written by unit TR-6 (plans/contraste/manuscripts-plan.md in CAOS_MANAGE).
% The tables and the numbers of this chapter are generated by tables/make_market.py from data/derived (cases C11, C13 and the size and
% power studies of C22) and the figures by figures/make_market.py; the \mk... macros come from tables/market-numbers.tex.
% Generated by tables/make_market.py from data/derived (C11, C13, C22) and docs/design/features/c11: do not edit.
% One macro per number the market section quotes.
\newcommand{\mkAboveCalm}{1}
\newcommand{\mkAboveCovid}{13}
\newcommand{\mkAboveFx}{7}
\newcommand{\mkAboveGfc}{18}
\newcommand{\mkAboveMixed}{0}
\newcommand{\mkAboveRates}{8}
\newcommand{\mkAmbHiMissingBasis}{97.0}
\newcommand{\mkAmbHiNone}{0.4}
\newcommand{\mkAmbHiOmittedFactor}{91.8}
\newcommand{\mkAmbHiStaleVolatility}{0.4}
\newcommand{\mkAmbHiTimeMisalignment}{100.0}
\newcommand{\mkAmbHiWrongCorrelation}{100.0}
\newcommand{\mkAmbLoMissingBasis}{94.6}
\newcommand{\mkAmbLoNone}{0.0}
\newcommand{\mkAmbLoOmittedFactor}{88.1}
\newcommand{\mkAmbLoStaleVolatility}{0.0}
\newcommand{\mkAmbLoTimeMisalignment}{99.6}
\newcommand{\mkAmbLoWrongCorrelation}{99.6}
\newcommand{\mkAmbMissingBasis}{96.0}
\newcommand{\mkAmbNone}{0.0}
\newcommand{\mkAmbOmittedFactor}{90.1}
\newcommand{\mkAmbStaleVolatility}{0.0}
\newcommand{\mkAmbTimeMisalignment}{100.0}
\newcommand{\mkAmbWrongCorrelation}{100.0}
\newcommand{\mkAmberAlpha}{0.1}
\newcommand{\mkAwLam}{0.99}
\newcommand{\mkAzAmber}{0.05}
\newcommand{\mkAzRed}{0.0001}
\newcommand{\mkBandHiCalm}{13}
\newcommand{\mkBandHiCovid}{3}
\newcommand{\mkBandHiFx}{6}
\newcommand{\mkBandHiGfc}{10}
\newcommand{\mkBandHiMixed}{3}
\newcommand{\mkBandHiRates}{6}
\newcommand{\mkBandLoCalm}{3}
\newcommand{\mkBandLoCovid}{0}
\newcommand{\mkBandLoFx}{0}
\newcommand{\mkBandLoGfc}{1}
\newcommand{\mkBandLoMixed}{0}
\newcommand{\mkBandLoRates}{0}
\newcommand{\mkBestAnyACalm}{GJR-t}
\newcommand{\mkBestAnyACovid}{Filtered HS}
\newcommand{\mkBestAnyAFx}{ES-CAViaR}
\newcommand{\mkBestAnyAGfc}{GJR-t}
\newcommand{\mkBestAnyAMixed}{GJR-t}
\newcommand{\mkBestAnyARates}{EWMA}
\newcommand{\mkBestAnyBCalm}{EWMA}
\newcommand{\mkBestAnyBCovid}{Quantile network}
\newcommand{\mkBestAnyBFx}{Cond. EVT}
\newcommand{\mkBestAnyBGfc}{GAS-FZ 1F}
\newcommand{\mkBestAnyBMixed}{GJR-t}
\newcommand{\mkBestAnyBRates}{EWMA}
\newcommand{\mkCRep}{1000}
\newcommand{\mkCWin}{250}
\newcommand{\mkCaviarRefit}{5}
\newcommand{\mkCaviarStarts}{20}
\newcommand{\mkChampACalm}{GJR-t}
\newcommand{\mkChampACovid}{Filtered HS}
\newcommand{\mkChampAFx}{ES-CAViaR}
\newcommand{\mkChampAGfc}{none}
\newcommand{\mkChampAMixed}{GJR-t}
\newcommand{\mkChampARates}{EWMA}
\newcommand{\mkChampBCalm}{EWMA}
\newcommand{\mkChampBCovid}{Cond. EVT}
\newcommand{\mkChampBFx}{Cond. EVT}
\newcommand{\mkChampBGfc}{none}
\newcommand{\mkChampBMixed}{GJR-t}
\newcommand{\mkChampBRates}{EWMA}
\newcommand{\mkChampFzACalm}{-4.521}
\newcommand{\mkChampFzACovid}{-1.986}
\newcommand{\mkChampFzAFx}{-3.557}
\newcommand{\mkChampFzAMixed}{-4.439}
\newcommand{\mkChampFzARates}{-3.811}
\newcommand{\mkChampFzBCalm}{-4.659}
\newcommand{\mkChampFzBCovid}{-2.364}
\newcommand{\mkChampFzBFx}{-3.716}
\newcommand{\mkChampFzBMixed}{-4.623}
\newcommand{\mkChampFzBRates}{-4.002}
\newcommand{\mkDaysCalm}{750}
\newcommand{\mkDaysCovid}{104}
\newcommand{\mkDaysFx}{251}
\newcommand{\mkDaysGfc}{504}
\newcommand{\mkDaysMixed}{64}
\newcommand{\mkDaysRates}{249}
\newcommand{\mkDeskExcNFMissingBasis}{22.2}
\newcommand{\mkDeskExcNFNone}{6.5}
\newcommand{\mkDeskExcNFOmittedFactor}{14.9}
\newcommand{\mkDeskExcNFStaleVolatility}{9.2}
\newcommand{\mkDeskExcNFTimeMisalignment}{6.5}
\newcommand{\mkDeskExcNFWrongCorrelation}{28.8}
\newcommand{\mkDeskExcNNMissingBasis}{14.2}
\newcommand{\mkDeskExcNNNone}{2.6}
\newcommand{\mkDeskExcNNOmittedFactor}{8.0}
\newcommand{\mkDeskExcNNStaleVolatility}{5.0}
\newcommand{\mkDeskExcNNTimeMisalignment}{2.7}
\newcommand{\mkDeskExcNNWrongCorrelation}{19.3}
\newcommand{\mkDeskFlagMissingBasis}{56.5}
\newcommand{\mkDeskFlagNone}{0.0}
\newcommand{\mkDeskFlagOmittedFactor}{6.4}
\newcommand{\mkDeskFlagStaleVolatility}{0.5}
\newcommand{\mkDeskFlagTimeMisalignment}{0.0}
\newcommand{\mkDeskFlagWrongCorrelation}{95.4}
\newcommand{\mkDeskLimitNF}{30}
\newcommand{\mkDeskLimitNN}{12}
\newcommand{\mkDnGapACalm}{0.14}
\newcommand{\mkDnGapACovid}{15.28}
\newcommand{\mkDnGapAFx}{2.68}
\newcommand{\mkDnGapAGfc}{4.91}
\newcommand{\mkDnGapAMixed}{0.40}
\newcommand{\mkDnGapARates}{0.81}
\newcommand{\mkDnGapBCalm}{0.13}
\newcommand{\mkDnGapBCovid}{7.38}
\newcommand{\mkDnGapBFx}{1.50}
\newcommand{\mkDnGapBGfc}{2.49}
\newcommand{\mkDnGapBMixed}{0.42}
\newcommand{\mkDnGapBRates}{0.45}
\newcommand{\mkEligCalm}{11}
\newcommand{\mkEligCovid}{4}
\newcommand{\mkEligFx}{11}
\newcommand{\mkEligGfc}{0}
\newcommand{\mkEligMixed}{18}
\newcommand{\mkEligRates}{4}
\newcommand{\mkEndCalm}{2006-12-29}
\newcommand{\mkEndCovid}{2020-06-30}
\newcommand{\mkEndFx}{2009-06-30}
\newcommand{\mkEndGfc}{2009-06-30}
\newcommand{\mkEndMixed}{2023-05-31}
\newcommand{\mkEndRates}{2022-12-30}
\newcommand{\mkEstWin}{1000}
\newcommand{\mkEvAzCalm}{30}
\newcommand{\mkEvAzCovid}{30}
\newcommand{\mkEvAzFx}{30}
\newcommand{\mkEvAzGfc}{30}
\newcommand{\mkEvAzMixed}{30}
\newcommand{\mkEvAzRates}{30}
\newcommand{\mkEvCcCalm}{35}
\newcommand{\mkEvCcCovid}{36}
\newcommand{\mkEvCcFx}{36}
\newcommand{\mkEvCcGfc}{36}
\newcommand{\mkEvCcMixed}{4}
\newcommand{\mkEvCcRates}{36}
\newcommand{\mkEvDqCalm}{36}
\newcommand{\mkEvDqCovid}{36}
\newcommand{\mkEvDqFx}{36}
\newcommand{\mkEvDqGfc}{36}
\newcommand{\mkEvDqMixed}{36}
\newcommand{\mkEvDqRates}{36}
\newcommand{\mkEvDurCalm}{34}
\newcommand{\mkEvDurCovid}{36}
\newcommand{\mkEvDurFx}{36}
\newcommand{\mkEvDurGfc}{36}
\newcommand{\mkEvDurMixed}{0}
\newcommand{\mkEvDurRates}{36}
\newcommand{\mkEvIndCalm}{35}
\newcommand{\mkEvIndCovid}{36}
\newcommand{\mkEvIndFx}{36}
\newcommand{\mkEvIndGfc}{36}
\newcommand{\mkEvIndMixed}{4}
\newcommand{\mkEvIndRates}{36}
\newcommand{\mkEvKupCalm}{36}
\newcommand{\mkEvKupCovid}{36}
\newcommand{\mkEvKupFx}{36}
\newcommand{\mkEvKupGfc}{36}
\newcommand{\mkEvKupMixed}{36}
\newcommand{\mkEvKupRates}{36}
\newcommand{\mkEvtFrac}{0.05}
\newcommand{\mkEwmaLam}{0.97}
\newcommand{\mkExcBDnCalm}{2}
\newcommand{\mkExcBDnCovid}{20}
\newcommand{\mkExcBDnFx}{35}
\newcommand{\mkExcBDnGfc}{77}
\newcommand{\mkExcBDnMixed}{0}
\newcommand{\mkExcBDnRates}{25}
\newcommand{\mkExcBEvtCalm}{15}
\newcommand{\mkExcBEvtCovid}{4}
\newcommand{\mkExcBEvtFx}{7}
\newcommand{\mkExcBEvtGfc}{26}
\newcommand{\mkExcBEvtMixed}{0}
\newcommand{\mkExcBEvtRates}{13}
\newcommand{\mkExcBEwmaCalm}{17}
\newcommand{\mkExcBEwmaCovid}{9}
\newcommand{\mkExcBEwmaFx}{8}
\newcommand{\mkExcBEwmaGfc}{26}
\newcommand{\mkExcBEwmaMixed}{0}
\newcommand{\mkExcBEwmaRates}{9}
\newcommand{\mkExcBFhsCalm}{15}
\newcommand{\mkExcBFhsCovid}{5}
\newcommand{\mkExcBFhsFx}{7}
\newcommand{\mkExcBFhsGfc}{24}
\newcommand{\mkExcBFhsMixed}{0}
\newcommand{\mkExcBFhsRates}{12}
\newcommand{\mkExcBGasCalm}{9}
\newcommand{\mkExcBGasCovid}{8}
\newcommand{\mkExcBGasFx}{21}
\newcommand{\mkExcBGasGfc}{44}
\newcommand{\mkExcBGasMixed}{0}
\newcommand{\mkExcBGasRates}{15}
\newcommand{\mkExcBGjrCalm}{13}
\newcommand{\mkExcBGjrCovid}{5}
\newcommand{\mkExcBGjrFx}{8}
\newcommand{\mkExcBGjrGfc}{30}
\newcommand{\mkExcBGjrMixed}{1}
\newcommand{\mkExcBGjrRates}{18}
\newcommand{\mkExcBHsCalm}{22}
\newcommand{\mkExcBHsCovid}{13}
\newcommand{\mkExcBHsFx}{13}
\newcommand{\mkExcBHsGfc}{36}
\newcommand{\mkExcBHsMixed}{0}
\newcommand{\mkExcBHsRates}{14}
\newcommand{\mkExcBNnCalm}{35}
\newcommand{\mkExcBNnCovid}{7}
\newcommand{\mkExcBNnFx}{10}
\newcommand{\mkExcBNnGfc}{23}
\newcommand{\mkExcBNnMixed}{1}
\newcommand{\mkExcBNnRates}{12}
\newcommand{\mkExcDnCalm}{0}
\newcommand{\mkExcDnCovid}{16}
\newcommand{\mkExcDnFx}{23}
\newcommand{\mkExcDnGfc}{57}
\newcommand{\mkExcDnMixed}{0}
\newcommand{\mkExcDnRates}{15}
\newcommand{\mkExcEvtCalm}{8}
\newcommand{\mkExcEvtCovid}{2}
\newcommand{\mkExcEvtFx}{3}
\newcommand{\mkExcEvtGfc}{13}
\newcommand{\mkExcEvtMixed}{0}
\newcommand{\mkExcEvtRates}{5}
\newcommand{\mkExcEwmaCalm}{10}
\newcommand{\mkExcEwmaCovid}{7}
\newcommand{\mkExcEwmaFx}{3}
\newcommand{\mkExcEwmaGfc}{16}
\newcommand{\mkExcEwmaMixed}{0}
\newcommand{\mkExcEwmaRates}{6}
\newcommand{\mkExcFhsCalm}{8}
\newcommand{\mkExcFhsCovid}{2}
\newcommand{\mkExcFhsFx}{3}
\newcommand{\mkExcFhsGfc}{13}
\newcommand{\mkExcFhsMixed}{0}
\newcommand{\mkExcFhsRates}{5}
\newcommand{\mkExcGasCalm}{3}
\newcommand{\mkExcGasCovid}{7}
\newcommand{\mkExcGasFx}{19}
\newcommand{\mkExcGasGfc}{26}
\newcommand{\mkExcGasMixed}{0}
\newcommand{\mkExcGasRates}{8}
\newcommand{\mkExcGjrCalm}{9}
\newcommand{\mkExcGjrCovid}{3}
\newcommand{\mkExcGjrFx}{3}
\newcommand{\mkExcGjrGfc}{14}
\newcommand{\mkExcGjrMixed}{0}
\newcommand{\mkExcGjrRates}{6}
\newcommand{\mkExcHsCalm}{9}
\newcommand{\mkExcHsCovid}{8}
\newcommand{\mkExcHsFx}{10}
\newcommand{\mkExcHsGfc}{19}
\newcommand{\mkExcHsMixed}{0}
\newcommand{\mkExcHsRates}{8}
\newcommand{\mkExcNnCalm}{22}
\newcommand{\mkExcNnCovid}{4}
\newcommand{\mkExcNnFx}{5}
\newcommand{\mkExcNnGfc}{16}
\newcommand{\mkExcNnMixed}{1}
\newcommand{\mkExcNnRates}{11}
\newcommand{\mkExpBCalm}{18.75}
\newcommand{\mkExpBCovid}{2.60}
\newcommand{\mkExpBFx}{6.28}
\newcommand{\mkExpBGfc}{12.60}
\newcommand{\mkExpBMixed}{1.60}
\newcommand{\mkExpBRates}{6.22}
\newcommand{\mkExpCalm}{7.50}
\newcommand{\mkExpCovid}{1.04}
\newcommand{\mkExpFx}{2.51}
\newcommand{\mkExpGfc}{5.04}
\newcommand{\mkExpMixed}{0.64}
\newcommand{\mkExpRates}{2.49}
\newcommand{\mkGBurn}{300}
\newcommand{\mkGLoadA}{0.006}
\newcommand{\mkGLoadB}{-0.004}
\newcommand{\mkGPhi}{0.98}
\newcommand{\mkGResShare}{0.1}
\newcommand{\mkGRho}{0.6}
\newcommand{\mkGVolSd}{0.35}
\newcommand{\mkGasRefit}{20}
\newcommand{\mkGasStarts}{20}
\newcommand{\mkGfcKupP}{0.00000171}
\newcommand{\mkGfcKupStat}{22.90}
\newcommand{\mkGjrRoughACalm}{0.026}
\newcommand{\mkGjrRoughACovid}{0.171}
\newcommand{\mkGjrRoughAFx}{0.020}
\newcommand{\mkGjrRoughAGfc}{0.058}
\newcommand{\mkGjrRoughAMixed}{0.058}
\newcommand{\mkGjrRoughARates}{0.066}
\newcommand{\mkHacLags}{0}
\newcommand{\mkHsLong}{500}
\newcommand{\mkHsRoughACalm}{0.002}
\newcommand{\mkHsRoughACovid}{0.012}
\newcommand{\mkHsRoughAFx}{0.002}
\newcommand{\mkHsRoughAGfc}{0.003}
\newcommand{\mkHsRoughAMixed}{0.002}
\newcommand{\mkHsRoughARates}{0.003}
\newcommand{\mkHsShort}{250}
\newcommand{\mkKsFlagMissingBasis}{86.7}
\newcommand{\mkKsFlagNone}{0.0}
\newcommand{\mkKsFlagOmittedFactor}{44.8}
\newcommand{\mkKsFlagStaleVolatility}{0.0}
\newcommand{\mkKsFlagTimeMisalignment}{0.0}
\newcommand{\mkKsFlagWrongCorrelation}{100.0}
\newcommand{\mkKsMissingBasis}{0.120}
\newcommand{\mkKsNone}{0.033}
\newcommand{\mkKsOmittedFactor}{0.089}
\newcommand{\mkKsSdMissingBasis}{0.028}
\newcommand{\mkKsSdNone}{0.006}
\newcommand{\mkKsSdOmittedFactor}{0.016}
\newcommand{\mkKsSdStaleVolatility}{0.008}
\newcommand{\mkKsSdTimeMisalignment}{0.006}
\newcommand{\mkKsSdWrongCorrelation}{0.014}
\newcommand{\mkKsStaleVolatility}{0.041}
\newcommand{\mkKsTimeMisalignment}{0.033}
\newcommand{\mkKsWrongCorrelation}{0.141}
\newcommand{\mkMcDraws}{10{,}000}
\newcommand{\mkMcsACalm}{13}
\newcommand{\mkMcsACovid}{15}
\newcommand{\mkMcsAFx}{14}
\newcommand{\mkMcsAGfc}{13}
\newcommand{\mkMcsAMixed}{2}
\newcommand{\mkMcsARates}{15}
\newcommand{\mkMcsBCalm}{13}
\newcommand{\mkMcsBCovid}{14}
\newcommand{\mkMcsBFx}{14}
\newcommand{\mkMcsBGfc}{13}
\newcommand{\mkMcsBMixed}{4}
\newcommand{\mkMcsBRates}{13}
\newcommand{\mkMcsBlock}{10}
\newcommand{\mkMcsBoot}{1000}
\newcommand{\mkMcsConf}{90}
\newcommand{\mkMcsMaxA}{15}
\newcommand{\mkMcsMinA}{13}
\newcommand{\mkMcsNACalm}{15}
\newcommand{\mkMcsNACovid}{15}
\newcommand{\mkMcsNAFx}{15}
\newcommand{\mkMcsNAGfc}{15}
\newcommand{\mkMcsNAMixed}{15}
\newcommand{\mkMcsNARates}{15}
\newcommand{\mkMcsNBCalm}{15}
\newcommand{\mkMcsNBCovid}{15}
\newcommand{\mkMcsNBFx}{15}
\newcommand{\mkMcsNBGfc}{15}
\newcommand{\mkMcsNBMixed}{15}
\newcommand{\mkMcsNBRates}{15}
\newcommand{\mkMissDaysCalm}{78}
\newcommand{\mkMissDaysCovid}{0}
\newcommand{\mkMissDaysFx}{43}
\newcommand{\mkMissDaysGfc}{0}
\newcommand{\mkMissDaysMixed}{50}
\newcommand{\mkMissDaysRates}{0}
\newcommand{\mkNGreenMissingBasis}{40}
\newcommand{\mkNGreenNone}{1000}
\newcommand{\mkNGreenOmittedFactor}{99}
\newcommand{\mkNGreenStaleVolatility}{1000}
\newcommand{\mkNGreenTimeMisalignment}{0}
\newcommand{\mkNGreenWrongCorrelation}{0}
\newcommand{\mkNLadder}{17}
\newcommand{\mkNModels}{18}
\newcommand{\mkNRankCalm}{15}
\newcommand{\mkNRankCovid}{15}
\newcommand{\mkNRankFx}{15}
\newcommand{\mkNRankGfc}{15}
\newcommand{\mkNRankMixed}{15}
\newcommand{\mkNRankRates}{15}
\newcommand{\mkNRedMissingBasis}{575}
\newcommand{\mkNRedNone}{0}
\newcommand{\mkNRedOmittedFactor}{43}
\newcommand{\mkNRedStaleVolatility}{0}
\newcommand{\mkNRedTimeMisalignment}{1000}
\newcommand{\mkNRedWrongCorrelation}{913}
\newcommand{\mkNStudies}{12}
\newcommand{\mkNWindows}{6}
\newcommand{\mkNnChampA}{0}
\newcommand{\mkNnChampB}{0}
\newcommand{\mkNnCost}{564{,}000}
\newcommand{\mkNnEligCalm}{no}
\newcommand{\mkNnEligCovid}{no}
\newcommand{\mkNnEligFx}{yes}
\newcommand{\mkNnEligGfc}{no}
\newcommand{\mkNnEligMixed}{yes}
\newcommand{\mkNnEligRates}{no}
\newcommand{\mkNnEligWindows}{2}
\newcommand{\mkNnEpochs}{200}
\newcommand{\mkNnExcBCalm}{35}
\newcommand{\mkNnExcBCovid}{7}
\newcommand{\mkNnExcBFx}{10}
\newcommand{\mkNnExcBGfc}{23}
\newcommand{\mkNnExcBMixed}{1}
\newcommand{\mkNnExcBRates}{12}
\newcommand{\mkNnExcCalm}{22}
\newcommand{\mkNnExcCovid}{4}
\newcommand{\mkNnExcFx}{5}
\newcommand{\mkNnExcGfc}{16}
\newcommand{\mkNnExcMixed}{1}
\newcommand{\mkNnExcRates}{11}
\newcommand{\mkNnFailures}{0}
\newcommand{\mkNnFzACalm}{-4.286}
\newcommand{\mkNnFzACovid}{-1.669}
\newcommand{\mkNnFzAFx}{-3.527}
\newcommand{\mkNnFzAGfc}{-2.371}
\newcommand{\mkNnFzAMixed}{-4.134}
\newcommand{\mkNnFzARates}{-3.445}
\newcommand{\mkNnFzBCalm}{-4.588}
\newcommand{\mkNnFzBCovid}{-2.368}
\newcommand{\mkNnFzBFx}{-3.685}
\newcommand{\mkNnFzBGfc}{-2.856}
\newcommand{\mkNnFzBMixed}{-4.374}
\newcommand{\mkNnFzBRates}{-3.909}
\newcommand{\mkNnHidden}{4}
\newcommand{\mkNnKupCalm}{$<0.001$}
\newcommand{\mkNnKupCovid}{0.026}
\newcommand{\mkNnKupFx}{0.164}
\newcommand{\mkNnKupGfc}{$<0.001$}
\newcommand{\mkNnKupMixed}{0.676}
\newcommand{\mkNnKupRates}{$<0.001$}
\newcommand{\mkNnLr}{0.02}
\newcommand{\mkNnParams}{58}
\newcommand{\mkNnRankCalm}{15}
\newcommand{\mkNnRankCovid}{4}
\newcommand{\mkNnRankFx}{2}
\newcommand{\mkNnRankGfc}{9}
\newcommand{\mkNnRankMixed}{9}
\newcommand{\mkNnRankRates}{11}
\newcommand{\mkNnRefit}{20}
\newcommand{\mkNnRefits}{200}
\newcommand{\mkNnRoughACalm}{0.161}
\newcommand{\mkNnRoughACovid}{0.373}
\newcommand{\mkNnRoughAFx}{0.211}
\newcommand{\mkNnRoughAGfc}{0.220}
\newcommand{\mkNnRoughAMixed}{0.241}
\newcommand{\mkNnRoughARates}{0.251}
\newcommand{\mkNnSetA}{5}
\newcommand{\mkNnSetB}{5}
\newcommand{\mkNnStarts}{3}
\newcommand{\mkNnTailA}{9.4}
\newcommand{\mkNnTailB}{23.5}
\newcommand{\mkNnTrainRows}{940}
\newcommand{\mkNnWd}{0.1}
\newcommand{\mkNnWin}{1000}
\newcommand{\mkOverNFMissingBasis}{13.0}
\newcommand{\mkOverNFNone}{0.0}
\newcommand{\mkOverNFOmittedFactor}{0.0}
\newcommand{\mkOverNFStaleVolatility}{0.0}
\newcommand{\mkOverNFTimeMisalignment}{0.0}
\newcommand{\mkOverNFWrongCorrelation}{36.7}
\newcommand{\mkOverNNMissingBasis}{56.5}
\newcommand{\mkOverNNNone}{0.0}
\newcommand{\mkOverNNOmittedFactor}{6.4}
\newcommand{\mkOverNNStaleVolatility}{0.5}
\newcommand{\mkOverNNTimeMisalignment}{0.0}
\newcommand{\mkOverNNWrongCorrelation}{95.3}
\newcommand{\mkPearsonMissingBasis}{0.756}
\newcommand{\mkPearsonNone}{0.995}
\newcommand{\mkPearsonOmittedFactor}{0.796}
\newcommand{\mkPearsonStaleVolatility}{0.955}
\newcommand{\mkPearsonTimeMisalignment}{0.000}
\newcommand{\mkPearsonWrongCorrelation}{0.990}
\newcommand{\mkPowAz}{89.4}
\newcommand{\mkPowDqLong}{27.9}
\newcommand{\mkPowDqShort}{16.0}
\newcommand{\mkPowInd}{87.5}
\newcommand{\mkPowKup}{69.3}
\newcommand{\mkRedAlpha}{0.05}
\newcommand{\mkRedAzCalm}{0}
\newcommand{\mkRedAzCovid}{14}
\newcommand{\mkRedAzFx}{9}
\newcommand{\mkRedAzGfc}{16}
\newcommand{\mkRedAzMixed}{0}
\newcommand{\mkRedAzRates}{4}
\newcommand{\mkRedCcCalm}{7}
\newcommand{\mkRedCcCovid}{21}
\newcommand{\mkRedCcFx}{13}
\newcommand{\mkRedCcGfc}{36}
\newcommand{\mkRedCcMixed}{0}
\newcommand{\mkRedCcRates}{18}
\newcommand{\mkRedDqCalm}{4}
\newcommand{\mkRedDqCovid}{24}
\newcommand{\mkRedDqFx}{20}
\newcommand{\mkRedDqGfc}{36}
\newcommand{\mkRedDqMixed}{0}
\newcommand{\mkRedDqRates}{25}
\newcommand{\mkRedDurCalm}{2}
\newcommand{\mkRedDurCovid}{2}
\newcommand{\mkRedDurFx}{1}
\newcommand{\mkRedDurGfc}{7}
\newcommand{\mkRedDurMixed}{0}
\newcommand{\mkRedDurRates}{4}
\newcommand{\mkRedIndCalm}{0}
\newcommand{\mkRedIndCovid}{0}
\newcommand{\mkRedIndFx}{3}
\newcommand{\mkRedIndGfc}{4}
\newcommand{\mkRedIndMixed}{0}
\newcommand{\mkRedIndRates}{0}
\newcommand{\mkRedKupCalm}{10}
\newcommand{\mkRedKupCovid}{25}
\newcommand{\mkRedKupFx}{13}
\newcommand{\mkRedKupGfc}{36}
\newcommand{\mkRedKupMixed}{0}
\newcommand{\mkRedKupRates}{21}
\newcommand{\mkRedMissingBasis}{57.5}
\newcommand{\mkRedNone}{0.0}
\newcommand{\mkRedOmittedFactor}{4.3}
\newcommand{\mkRedStaleVolatility}{0.0}
\newcommand{\mkRedTimeMisalignment}{100.0}
\newcommand{\mkRedWrongCorrelation}{91.3}
\newcommand{\mkRoughMaxEwma}{0.030}
\newcommand{\mkRoughMaxGjr}{0.171}
\newcommand{\mkRoughMaxGjrOut}{0.066}
\newcommand{\mkRoughMaxHs}{0.012}
\newcommand{\mkRoughMaxNn}{0.373}
\newcommand{\mkRoughMinEwma}{0.011}
\newcommand{\mkRoughMinGjr}{0.020}
\newcommand{\mkRoughMinHs}{0.002}
\newcommand{\mkRoughMinNn}{0.161}
\newcommand{\mkSeed}{42}
\newcommand{\mkSelAgeWeightLambda}{0.99}
\newcommand{\mkSelEstimationWindowDays}{1000}
\newcommand{\mkSelEvtThresholdFraction}{0.05}
\newcommand{\mkSelEwmaLambda}{0.97}
\newcommand{\mkSelNnHiddenUnits}{4}
\newcommand{\mkSelNnWeightDecay}{0.1}
\newcommand{\mkSelWinGap}{0.08}
\newcommand{\mkSignif}{5}
\newcommand{\mkSizeAboveFive}{2}
\newcommand{\mkSizeEdge}{1}
\newcommand{\mkSizeMax}{6.7}
\newcommand{\mkSizeMin}{4.2}
\newcommand{\mkSizeOverCc}{6.7}
\newcommand{\mkSizeOverCcHi}{7.8}
\newcommand{\mkSizeOverCcLo}{5.6}
\newcommand{\mkSizeOverInd}{6.6}
\newcommand{\mkSizeOverIndHi}{7.8}
\newcommand{\mkSizeOverIndLo}{5.6}
\newcommand{\mkSpFlagMissingBasis}{91.9}
\newcommand{\mkSpFlagNone}{0.0}
\newcommand{\mkSpFlagOmittedFactor}{80.5}
\newcommand{\mkSpFlagStaleVolatility}{0.0}
\newcommand{\mkSpFlagTimeMisalignment}{100.0}
\newcommand{\mkSpFlagWrongCorrelation}{0.0}
\newcommand{\mkSpMissingBasis}{0.705}
\newcommand{\mkSpNone}{0.992}
\newcommand{\mkSpOmittedFactor}{0.772}
\newcommand{\mkSpSdMissingBasis}{0.073}
\newcommand{\mkSpSdNone}{0.004}
\newcommand{\mkSpSdOmittedFactor}{0.031}
\newcommand{\mkSpSdStaleVolatility}{0.005}
\newcommand{\mkSpSdTimeMisalignment}{0.069}
\newcommand{\mkSpSdWrongCorrelation}{0.004}
\newcommand{\mkSpStaleVolatility}{0.979}
\newcommand{\mkSpTimeMisalignment}{-0.004}
\newcommand{\mkSpWrongCorrelation}{0.986}
\newcommand{\mkSpanEnd}{2003-12-31}
\newcommand{\mkSpanStart}{1995-01-01}
\newcommand{\mkSpanStride}{5}
\newcommand{\mkStartCalm}{2004-01-02}
\newcommand{\mkStartCovid}{2020-02-03}
\newcommand{\mkStartFx}{2008-07-01}
\newcommand{\mkStartGfc}{2007-07-02}
\newcommand{\mkStartMixed}{2023-03-01}
\newcommand{\mkStartRates}{2022-01-03}
\newcommand{\mkThrKsGreen}{0.09}
\newcommand{\mkThrKsRed}{0.12}
\newcommand{\mkThrSpGreen}{0.8}
\newcommand{\mkThrSpRed}{0.7}
\newcommand{\mkTlAmber}{60}
\newcommand{\mkTlFinalRed}{6}
\newcommand{\mkTlFiveCum}{95.88}
\newcommand{\mkTlFiveMult}{1.70}
\newcommand{\mkTlFiveUs}{3.40}
\newcommand{\mkTlFourCum}{89.22}
\newcommand{\mkTlFourMult}{1.50}
\newcommand{\mkTlFourUs}{3.00}
\newcommand{\mkTlGreen}{0}
\newcommand{\mkTlMax}{14}
\newcommand{\mkTlN}{255}
\newcommand{\mkTlNineCum}{99.97}
\newcommand{\mkTlNineMult}{1.92}
\newcommand{\mkTlNineUs}{3.85}
\newcommand{\mkTlReason}{fewer than 250 forecast days}
\newcommand{\mkTlRed}{195}
\newcommand{\mkTlTenCum}{99.99}
\newcommand{\mkTlTenMult}{2.00}
\newcommand{\mkTlTenUs}{4.00}
\newcommand{\mkTotalDays}{1922}
\newcommand{\mkZDraws}{10{,}000}
\newcommand{\mkZRefAlpha}{2.5}
\newcommand{\mkZRefN}{250}
\newcommand{\mkZRefRed}{1.8}
\newcommand{\mkZRefYellow}{0.70}

\emergencystretch=3em
\providecommand{\mkfn}[1]{\nolinkurl{#1}}

\section{Methods: market risk: value at risk, expected shortfall and their back-tests}
\label{sec:market}

Market risk is the loss that a portfolio of traded positions suffers when prices, interest rates, exchange rates or volatilities move.
A bank summarises it, day by day, with two numbers at a tail probability $\alpha$: the value at risk (VaR), the return that is
exceeded on the downside with probability $\alpha$, and the expected shortfall (ES), the mean of the returns beyond it. With $y_t$ the
return of the portfolio on day $t$ as a decimal of its value and $F_t$ the distribution function of $y_t$ given everything observed
before day $t$, the pair of that day is
\begin{equation}
v_t=F_t^{-1}(\alpha),\qquad e_t=\frac1\alpha\int_0^\alpha F_t^{-1}(u)\dd u\;\le\;v_t,
\label{eq:mk-def}
\end{equation}
both negative numbers in practice. The engine of the product stores them in this return convention, which the scoring rule of
\cref{sec:m18} requires, and the documents state positive loss numbers $-v_t$ and $-e_t$; the sign is flipped once, at the reporting
layer, so that no page mixes the two conventions. The pair is preferred to either number alone for reasons that are mathematical
before they are regulatory. The VaR fails subadditivity in general, so the VaR of a merged book can exceed the sum of the VaRs of its
parts \citep{market:artzner1999}; the ES, defined through the tail integral so that atoms are handled, is a coherent measure
\citep{credit-capital:acerbitasche2002}. Quantiles are elicitable, which means that a scoring function exists whose expected value is minimised
by the true quantile, and the ES is not elicitable on its own \citep{market:gneiting2011}; the pair $(v_t,e_t)$ is, however, jointly
elicitable, and the full class of strictly consistent scoring functions for it is known \citep{market:fisslerziegel2016}. That
property lets one loss function both estimate and rank the two numbers (\cref{sec:m18}) and is what makes a comparison of ES forecasts
possible at all.

The regulation that governs the area is the supervisory back-testing of internal models. The Basel Committee introduced it in 1996:
the days on which the loss exceeded the 99 percent VaR in the last 250 are counted, and the count places the model in a green, a yellow
or a red zone \citep{market:bcbs1996}. The consolidated Basel Framework keeps the count and the zones in chapter MAR32, with a
multiplier on the capital that rises with the count, and adds two desk-level requirements: a desk back-test, under which a desk with
more than \mkDeskLimitNN{} exceptions of its 99 percent VaR or more than \mkDeskLimitNF{} exceptions of its 97.5 percent VaR in twelve
months moves to the standardised approach, and the profit-and-loss attribution test, which compares the hypothetical P\&L of a desk with the
P\&L its risk model assigns to the same positions through a Spearman correlation and a Kolmogorov--Smirnov distance
\citep{market:bcbsmar32}. The capital measure of the internal models approach is the ES at 97.5 percent \citep{market:bcbsmar33}. The
multiplication factors of the United States rule, \mkTlFourUs{} to \mkTlTenUs{} instead of \mkTlFourMult{} to \mkTlTenMult{}, are the same table
with another scale \citep{market:cfr217204}. The legal texts say how a bank is judged; they do not say which model should be used, how long a window
has to be to say anything, or whether a test that passes a model says it is right, and those are the questions of a validator.

Two cases of the product put the questions to data. Case C11 is the market case on real outcomes. It forecasts, one day ahead and
strictly out of sample, the 99 and the 97.5 percent VaR and the ES of \mkNModels{} models (seventeen estimators and a quantile network)
through \mkNWindows{} windows that hold \mkTotalDays{} forecast days in all: the equity crisis of \mkStartGfc{} to \mkEndGfc{}, the
euro against the dollar in \mkStartFx{} to \mkEndFx{}, the pandemic crash of \mkStartCovid{} to \mkEndCovid{}, the hiking cycle of
2022, the bank turmoil of \mkStartMixed{} to \mkEndMixed{}, and, as the control, the calm years \mkStartCalm{} to \mkEndCalm{}. The
books are built from three public daily series: the return of the United States equity market, taken from the factors of the Kenneth
French data library; a long position in the 10-year Treasury par bond, whose return is minus its modified duration at the previous
day's yield times the yield change, with carry and convexity left out; and the euro in United States dollars of the Federal Reserve H.10
release. A window is one of these books or the mixed book of 50 percent equity, 30 percent bond and 20 percent euro, rebalanced daily.
The case asks which model is eligible, which model beats the champion under a consistent score, whether exceedances cluster, and how
the answers change between a calm year and a crisis. Case C13 is the desk case, on generated data with a known truth: a two-factor
spread desk whose risk model carries one planted deficiency at a time (a missing basis, a volatility input lagged by 20 days, a
correlation 0.3 too high, an omitted risk factor, a P\&L dated one day off) or none, in \mkCRep{} independent replications of
\mkCWin{} days each. It asks whether the attribution test and the desk back-test, the two statistical tests a desk must pass to stay on
internal models, catch each deficiency and how often.

Three methods carry the section. The estimators that turn a window of returns into the pair of equation~\eqref{eq:mk-def}, and the
scoring rule that ranks them, are the subject of \cref{sec:m18}. The back-tests that ask whether a forecast was right, the comparison
tests that ask which of two forecasts was better, the eligibility and champion rule of C11, and the two desk tests of C13, are the
subject of \cref{sec:m19}. A quantile network trained directly on the scoring rule, the one neural estimator in the ladder, is
developed in \cref{sec:m30}; the results of all three on all cases are in \cref{sec:results-market}.

% ----------------------------------------------------------------------------------------------------------------------------------
\subsection{Value-at-risk and expected shortfall estimators}
\label{sec:m18}

\subsubsection{The estimation problem}
For the forecast day $t$ an estimator receives the returns before $t$ and returns the pair $(\widehat v_t,\widehat e_t)$ of
equation~\eqref{eq:mk-def}. Nothing of day $t$ or later enters: the estimation window is the $N$ returns $y_{t-N},\dots,y_{t-1}$, the
parameters of a model that is refitted every $R$ forecast days come from an earlier day and never from a later one, and the filter
that carries them to day $t$ runs over the current window. The estimators fall in five groups that differ in what they assume about
the distribution of $y_t$: the empirical distribution of the window with or without weights (historical simulation); a parametric
distribution with a moving scale (delta-normal, Cornish--Fisher, the exponentially weighted moving average); a volatility model with a
parametric innovation law or an empirical one (the GARCH family, filtered historical simulation) or with a generalised Pareto tail
(conditional extreme value theory); a model of the quantile itself (CAViaR and its ES extension); and a score-driven model estimated
on the joint loss of the pair (the GAS-FZ models). The quantile network of \cref{sec:m30} is the sixth group. Each is developed here
with its equation, its assumptions and its source; the ladder that collects them, with the parameters in force in case C11, is
\cref{tab:market-ladder}.

\begin{table}[htbp]
\centering
\scriptsize
\setlength{\tabcolsep}{2.5pt}
\caption{The eighteen forecasting models of case C11 as the committed models document records them: the estimation window in days, the number of forecast days between two refits, the hyperparameters fixed on the span 1995 to 2003 (or by the compute limit for the network) and the rule that turns the fitted model into an expected shortfall.}
\label{tab:market-ladder}
\begin{tabular}{p{3.6cm}rrp{2.6cm}p{5.4cm}}
\toprule
Model & Window (days) & Refit (days) & Setting & Expected shortfall rule \\
\midrule
Historical simulation, 250 days & 250 & 1 & -- & Tail mean of the sorted window \\
Historical simulation, 500 days & 500 & 1 & -- & Tail mean of the sorted window \\
Age-weighted historical simulation & 250 & 1 & $\lambda$ 0.99 & Tail integral of the interpolated weighted quantile function \\
Delta-normal & 1000 & 1 & -- & Normal ES from the window mean and volatility \\
Cornish-Fisher & 1000 & 1 & -- & None: the Cornish-Fisher expansion gives the VaR only \\
EWMA & 1000 & 1 & $\lambda$ 0.97 & Normal ES from the EWMA volatility \\
GARCH(1,1)-t & 1000 & 1 & -- & Student-t tail mean of the unit-variance innovation \\
GJR-GARCH-t & 1000 & 1 & -- & Student-t tail mean of the unit-variance innovation \\
EGARCH-t & 1000 & 1 & -- & Student-t tail mean of the unit-variance innovation \\
Filtered historical simulation & 1000 & 1 & -- & Tail mean of the GARCH-filtered standardised residuals \\
Conditional EVT & 1000 & 1 & tail share 0.05 & GPD tail ES, undefined for xi at or above 1 \\
CAViaR, absolute value & 1000 & 5 & -- & None: CAViaR forecasts the VaR only \\
CAViaR, asymmetric slope & 1000 & 5 & -- & None: CAViaR forecasts the VaR only \\
ES-CAViaR (two-step) & 1000 & 5 & -- & Autoregressive gap below the CAViaR VaR, fitted by the asymmetric Laplace likelihood \\
GAS-FZ, one factor & 1000 & 20 & -- & Joint VaR and ES from one dynamic factor (FZ0 estimation) \\
GAS-FZ, two factors & 1000 & 20 & -- & Joint VaR and ES from two dynamic factors (FZ0 estimation) \\
Hybrid GARCH-FZ & 1000 & 20 & -- & One dynamic factor with a log absolute return term (FZ0 estimation) \\
Quantile network (FZ0) & 1000 & 20 & hidden units 4; epochs 200; learning rate 0.02; weight decay 0.1; starts 3 & Second network output: the gap of the ES below the VaR, trained jointly on FZ0 \\
\bottomrule
\end{tabular}
\end{table}


\subsubsection{Historical simulation, age-weighted and filtered}
Historical simulation reads the pair from the order statistics of the window \citep{market:pritsker2001}. With $r_{(1)}\le\dots\le r_{(N)}$
the sorted returns and $k=\lceil\alpha N\rceil$, the VaR is the left-continuous inverse of the empirical distribution function and the
ES the empirical tail integral,
\begin{equation}
\widehat v=r_{(k)},\qquad \widehat e=\frac{\sum_{i<k}r_{(i)}+(\alpha N-(k-1))\,r_{(k)}}{\alpha N}.
\label{eq:mk-sample}
\end{equation}
When $\alpha N$ is not an integer the $k$-th return carries the fractional weight $\alpha N-(k-1)$, as in the tail integral of
\citet{credit-capital:acerbitasche2002}; the mean of the $k$ smallest returns would overstate the tail count. The window of 250 days at the 99
percent level holds $2.5$ tail observations, which is the reason the estimator is slow to react to a volatility regime and keeps the
memory of a crisis day until it leaves the window \citep{market:pritsker2001}; the case runs windows of \mkHsShort{} and \mkHsLong{} days.

The age-weighted variant gives the return observed $i$ days ago the weight
\begin{equation}
w_i=\frac{1-\lambda}{1-\lambda^{K}}\,\lambda^{\,i-1},\qquad i=1,\dots,K,\qquad\sum_{i=1}^{K}w_i=1,
\label{eq:mk-awhs}
\end{equation}
sorts the returns, cumulates the weights and interpolates linearly between adjacent points \citep{market:boudoukh1998}. With $C_j$ the
cumulated weight of the sorted return $j$, the quantile function of the weighted sample has two knots per return,
$(C_j-w_j/2,\,r_j)$ and $(C_j,\,(r_j+r_{j+1})/2)$, because half of the weight of a return lies on each side of it; the VaR is the
value of this piecewise linear function at $\alpha$. The paper gives no rule for the ES. The implementation takes the tail integral
$\frac1\alpha\int_0^\alpha Q(u)\dd u$ of the same piecewise linear function, evaluated exactly with trapezoids, which is the
definition of equation~\eqref{eq:mk-def} applied to the distribution from which the VaR is read; the rule is a decision of the
implementation and not of the paper. The decay is $\lambda=\mkAwLam{}$ in the case, fixed on the selection span.

Filtered historical simulation removes the volatility clustering from the returns before the simulation. With a volatility filter
that gives the conditional mean $\mu_t$ and scale $\sigma_t$ of every day, the standardised residuals $z_t=(r_t-\mu_t)/\sigma_t$ are
resampled at the forecast scale \citep{market:baronevosper1999},
\begin{equation}
r^{*}_{T+1,s}=\mu_{T+1}+\sigma_{T+1}\,z_s,\qquad z_t=\frac{r_t-\mu_t}{\sigma_t},
\label{eq:mk-fhs}
\end{equation}
and $(\widehat v,\widehat e)$ is read from the scenarios $r^*_{T+1,s}$ by equation~\eqref{eq:mk-sample}. The implementation uses every
standardised residual once, without a random bootstrap, so that the forecast is deterministic. The volatility-updated variant of
\citet{market:hullwhite1998} rescales the raw returns, $r^*_i=r_i\,\widehat\sigma_{T+1}/\widehat\sigma_i$, and is in the engine but
not in the ladder of C11. Filtered historical simulation passes the misspecification of its filter through to the forecast; with a
constant volatility it reduces to the plain estimator of equation~\eqref{eq:mk-sample}.

\subsubsection{Parametric estimators: delta-normal, Cornish--Fisher and the exponentially weighted average}
The delta-normal estimator assumes that the return over the horizon is normal with mean $\mu$ and volatility $\sigma>0$, the
assumption of the RiskMetrics system \citep{market:riskmetrics1996}. The VaR is the normal quantile and the ES follows from the
definition by the substitution $z=\Phi^{-1}(u)$:
\begin{equation}
\widehat v=\mu+\Phi^{-1}(\alpha)\,\sigma,\qquad
\widehat e=\mu+\frac{\sigma}{\alpha}\int_{-\infty}^{\Phi^{-1}(\alpha)}z\,\varphi(z)\dd z=\mu-\sigma\,\frac{\varphi\!\left(\Phi^{-1}(\alpha)\right)}{\alpha},
\label{eq:mk-dn}
\end{equation}
where $\Phi$ and $\varphi$ are the standard normal distribution and density, because $z\varphi(z)=-\varphi'(z)$. The result is affine
in $\mu$ and linear in $\sigma$, and $\widehat e\le\widehat v$ always. The estimator fails on optionality and on heavy tails; it is
kept as the transparent baseline. In case C11 the moments are those of the last \mkEstWin{} returns.

The Cornish--Fisher estimator corrects the normal quantile for the skewness $S$ and the excess kurtosis $K_e$ of the window with the
expansion of \citet{market:cornishfisher1938},
\begin{equation}
z_{CF}=z+\frac{(z^2-1)S}{6}+\frac{(z^3-3z)K_e}{24}-\frac{(2z^3-5z)S^2}{36},\qquad z=\Phi^{-1}(\alpha),\qquad\widehat v=\mu+z_{CF}\,\sigma.
\label{eq:mk-cf}
\end{equation}
The expansion gives a VaR only; the model has no ES. It is not monotone in $z$ for large $S$ and $K_e$, so a forecast from a
non-monotone map would not be the quantile of any distribution. The implementation therefore reports a validity flag: the derivative
\begin{equation}
\frac{\mathrm d z_{CF}}{\mathrm d z}=1+\frac{2zS}{6}+\frac{(3z^2-3)K_e}{24}-\frac{(6z^2-5)S^2}{36}
\label{eq:mk-cf-deriv}
\end{equation}
must be positive at every point of $[\Phi^{-1}(\alpha),0]$, checked on a grid of 201 points. The criterion is a choice of the
implementation, and it certifies the map on the tail up to the quantile only.

The exponentially weighted moving average estimates the variance recursively, with the decay $\lambda$ \citep{market:riskmetrics1996},
\begin{equation}
\sigma^2_{t+1\mid t}=\lambda\,\sigma^2_{t\mid t-1}+(1-\lambda)\,r_t^2,\qquad K=\frac{\ln\epsilon}{\ln\lambda},
\label{eq:mk-ewma}
\end{equation}
where $K$ is the number of days after which an observation keeps a share $\epsilon$ of its initial weight. The mean is assumed to be
zero, which is part of the method. The decay is $\lambda=\mkEwmaLam{}$ in the case, for which $K\approx151$ days at $\epsilon=1$
percent; the VaR and the ES are those of the delta-normal estimator of equation~\eqref{eq:mk-dn} with the latest variance of the
recursion. The normal quantile under-covers at the 99 percent level, and the initial variance is part of the estimation window and
never of the forecast day.

\subsubsection{Volatility models and extreme value theory}
The GARCH family models the conditional variance of the shock $\varepsilon_t=r_t-\mu$. Three specifications are fitted by
quasi-maximum likelihood: the GARCH(1,1) of \citet{market:bollerslev1986}, the GJR extension of \citet{market:gjr1993}, which adds the
asymmetry $\gamma$ for a negative shock, and the exponential model of \citet{market:nelson1991}, which needs no positivity constraint,
\begin{align}
\sigma_t^2&=\omega+\alpha_1\varepsilon_{t-1}^2+\beta_1\sigma_{t-1}^2, \label{eq:mk-garch}\\
\sigma_t^2&=\omega+\left(\alpha_1+\gamma\,\indic\{\varepsilon_{t-1}<0\}\right)\varepsilon_{t-1}^2+\beta_1\sigma_{t-1}^2, \label{eq:mk-gjr}\\
\ln\sigma_t^2&=\omega+\beta\ln\sigma_{t-1}^2+\theta z_{t-1}+\gamma\left(|z_{t-1}|-\E|z_{t-1}|\right), \label{eq:mk-egarch}
\end{align}
with $z_t=\varepsilon_t/\sigma_t$ the standardised shock. The fit is accepted only when the optimiser converged and the parameters are in
their domain, which differs by model: for the GARCH, $\omega>0$, $\alpha_1\ge0$, $\beta_1\ge0$ and $\alpha_1+\beta_1<1$; for the GJR, $\omega>0$, $\alpha_1\ge0$, $\beta_1\ge0$,
$\alpha_1+\gamma\ge0$ (so that $\gamma$ may be negative) and $\alpha_1+\gamma/2+\beta_1<1$; for the exponential model only $|\beta|<1$; and a degrees-of-freedom parameter $\nu>2$ for the Student-$t$ innovation. A failed fit keeps the
previous one and is counted; no parameter is ever invented. With the filter run through the window and the one-step variance
$\sigma_{t+1}^2$ read off, the forecast is
\begin{equation}
\widehat v=\mu+\sigma_{t+1}\,q_\alpha,\qquad\widehat e=\mu+\sigma_{t+1}\,\frac1\alpha\int_0^\alpha q_u\dd u,
\label{eq:mk-garch-fc}
\end{equation}
with $q$ the quantile function of the unit-variance innovation. For the Student-$t$ innovation $q_\alpha=t_\nu^{-1}(\alpha)\sqrt{(\nu-2)/\nu}$
and the integral has the closed form
\begin{equation}
\frac1\alpha\int_0^\alpha q_u\dd u=-\frac{f_\nu(t_\alpha)}{\alpha}\,\frac{\nu+t_\alpha^{2}}{\nu-1}\sqrt{\frac{\nu-2}{\nu}},\qquad t_\alpha=t_\nu^{-1}(\alpha),
\label{eq:mk-t-es}
\end{equation}
where $f_\nu$ is the Student-$t$ density. The engine tests the closed form against the quadrature of the definition and the one-step variance against the
filter of the \texttt{arch} package with the same fixed parameters. The three models of case C11 use Student-$t$ innovations and a
constant mean, on a window of \mkEstWin{} returns refitted every day. The GARCH misses the leverage of equity returns, which the GJR
and the exponential models add; quasi-maximum likelihood with a heavy tail needs a long window, which is why the window length is a
selected hyperparameter (\cref{tab:market-selection}).

Conditional extreme value theory fits the GARCH filter by pseudo-maximum likelihood with normal innovations and a constant mean, then
models the tail of the standardised loss residuals $-z_t$ above a high threshold $u$ with a generalised Pareto distribution
\citep{market:mcneilfrey2000}. With $n$ residuals, $N_u$ of them above $u$, and the fitted scale $\widehat\beta$ and shape $\widehat\xi$,
\begin{equation}
\widehat z_q=u+\frac{\widehat\beta}{\widehat\xi}\left[\left(\frac n{N_u}(1-q)\right)^{-\widehat\xi}-1\right],\qquad
\widehat{\ES}_q=\frac{\widehat z_q}{1-\widehat\xi}+\frac{\widehat\beta-\widehat\xi u}{1-\widehat\xi}\quad(\widehat\xi<1),
\label{eq:mk-evt}
\end{equation}
and the forecast in the return convention is $\widehat v=\mu_{t+1}-\sigma_{t+1}\widehat z_q$ and $\widehat e=\mu_{t+1}-\sigma_{t+1}\widehat{\ES}_q$.
The threshold $u$ is the empirical $(1-\mkEvtFrac{})$ quantile of the residuals, the excesses are fitted by maximum likelihood with the
location fixed at zero, and a fit with fewer than 10 exceedances is refused. When $\widehat\xi\ge1$ the tail mean is infinite and the
ES is undefined: the day is recorded as such and no number is reported. The threshold is the main model risk; the shape is unstable
in short windows. The AR(1) mean term of the original paper is not used.

\subsubsection{Quantile regression and score-driven models}
The CAViaR models of \citet{market:engleManganelli2004} model the VaR $f_t$ directly as an autoregression on its own past and on the
last return. The symmetric absolute value and the asymmetric slope specifications are
\begin{align}
f_t&=b_1+b_2f_{t-1}+b_3|y_{t-1}|, \label{eq:mk-sav}\\
f_t&=b_1+b_2f_{t-1}+b_3\,(y_{t-1})^{+}+b_4\,(y_{t-1})^{-}, \label{eq:mk-as}
\end{align}
with $(y)^+=\max(y,0)$ and $(y)^-=-\min(y,0)$. The coefficients minimise the tick loss
\begin{equation}
\sum_t\left(\alpha-\indic\{y_t<f_t\}\right)(y_t-f_t),
\label{eq:mk-tick}
\end{equation}
whose expected value is minimised by the true quantile. The objective is not convex, so the fit uses \mkCaviarStarts{} starts, the
first of them the constant benchmark $b=(q,0,0)$ with $q$ the empirical quantile of the window, and polishes each with two
Nelder--Mead runs; the reported loss therefore never exceeds that of any constant. The recursion is a linear filter. CAViaR gives the VaR
only. The ES extension of \citet{market:taylor2019} writes $\ES_t=\mathrm{VaR}_t-x_t$ with a nonnegative gap that is updated only
after an exceedance,
\begin{equation}
x_t=\begin{cases}\gamma_0+\gamma_1\left(\mathrm{VaR}_{t-1}-r_{t-1}\right)+\gamma_2\,x_{t-1}&r_{t-1}\le\mathrm{VaR}_{t-1},\\x_{t-1}&\text{otherwise},\end{cases}
\label{eq:mk-escaviar}
\end{equation}
and is estimated by maximising the asymmetric Laplace likelihood \citep{market:taylor2019,market:liuluger2026}
\begin{equation}
f(r)=\frac{\alpha-1}{\ES_t}\exp\!\left(\frac{(r-\mathrm{VaR}_t)\left(\alpha-\indic\{r\le\mathrm{VaR}_t\}\right)}{\alpha\,\ES_t}\right).
\label{eq:mk-ald}
\end{equation}
The implementation is two-step: the VaR path of the tick-loss fit is fixed, and the likelihood is maximised over $(\gamma_0,\gamma_1,\gamma_2)$ with
$\gamma_0,\gamma_1\ge0$ and $0\le\gamma_2<1$, so that $x_t\ge0$ and $\ES_t\le\mathrm{VaR}_t$ hold on every day; joint estimation of the
paper, over the VaR and the ES parameters together, would be a different estimator, and the two-step choice is a declared deviation.

The GAS-FZ models of \citet{market:pattonziegelchen2019} drive the pair $(v_t,e_t)$ by the score of the joint loss of equation~\eqref{eq:mk-fz0}
below. The one-factor model ties the pair to a single dynamic factor and the hybrid model adds the logarithm of the absolute return to the
factor recursion,
\begin{align}
v_t&=a\,e^{\kappa_t},\quad e_t=b\,e^{\kappa_t},\quad b<a<0,\qquad
\kappa_{t+1}=\beta\kappa_t+\gamma\,\frac{-1}{e_t}\left(\frac1\alpha\indic\{y_t\le v_t\}\,y_t-e_t\right), \label{eq:mk-gas1}\\
\kappa_{t+1}&=\beta\kappa_t+\gamma\,\frac{-1}{e_t}\left(\frac1\alpha\indic\{y_t\le v_t\}\,y_t-e_t\right)+\delta\ln|y_t|. \label{eq:mk-gash}
\end{align}
The two-factor model updates the two numbers by their own scores, $(v_{t+1},e_{t+1})'=w+B(v_t,e_t)'+A(\lambda_{v,t},\lambda_{e,t})'$ with
$\lambda_{v,t}=-v_t(\indic\{y_t\le v_t\}-\alpha)$ and $\lambda_{e,t}=\frac1\alpha\indic\{y_t\le v_t\}y_t-e_t$, a diagonal matrix $B$ and a
full matrix $A$. The intercept of the one-factor and hybrid recursions is not identified under this loss and is fixed at zero. Estimation
minimises the average FZ0 loss of the window by Nelder--Mead from \mkGasStarts{} starts, the first of them the constant empirical VaR and
ES, so the fitted loss never exceeds that benchmark's; a path that leaves the domain $e_t\le v_t<0$ costs a large constant, so every fitted
path stays inside it. A fit that does not converge is a result and not an error. The recursions react to a return only through a VaR
violation, and the estimation is unstable for small $\alpha$ and short windows.

\subsubsection{The joint scoring rule}
All the estimators above are ranked by one loss. Let $Y$ be a realised return, $(v,e)$ a forecast pair with $e\le v<0$, and $G_1$,
$G_2$ functions with $G_1$ weakly increasing, $G_2$ strictly increasing and positive, $\mathcal G_2'=G_2$. The Fissler--Ziegel class of
strictly consistent losses for the pair \citep{market:fisslerziegel2016,market:pattonziegelchen2019} is
\begin{equation}
L_{FZ}(Y,v,e;\alpha)=\left(\indic\{Y\le v\}-\alpha\right)\left(G_1(v)-G_1(Y)\right)+G_2(e)\Big(e-v+\frac1\alpha\indic\{Y\le v\}(v-Y)\Big)-\mathcal G_2(e).
\label{eq:mk-fz}
\end{equation}
The choice $G_1=0$, $G_2(x)=-1/x$, $\mathcal G_2(x)=-\ln(-x)$ is the only one for which the differences of the loss between two forecasts
are homogeneous of degree zero in the scale of the data \citep{market:pattonziegelchen2019}; substituting it in
equation~\eqref{eq:mk-fz} gives the FZ0 loss,
\begin{equation}
L_{FZ0}(Y,v,e;\alpha)=-\frac1{\alpha e}\,\indic\{Y\le v\}\,(v-Y)+\frac ve+\ln(-e)-1.
\label{eq:mk-fz0}
\end{equation}
It has no tuning parameter, which makes it the default rule of the product; the loss is dimensionless, it can be negative, and only
differences between models are meaningful. A forecaster without an ES (the Cornish--Fisher model and the two CAViaR models) has a tick
loss, the loss of equation~\eqref{eq:mk-tick} at the single quantile, but no FZ0 value, never a number invented for it. The average FZ0 loss over
the forecast days of a window and the tests that compare it between models are the comparative evidence of \cref{sec:m19}.

\subsubsection{Parameters, complexity and implementation}
The parameters of the case are in \cref{tab:market-ladder}. Four of them are chosen on a span that ends before the earliest test
window, so that no window tunes on itself: the estimation window among 500, 750 and 1,000 days, the EWMA decay among 0.94, 0.97 and
0.99, the age-weighting decay among 0.98 and 0.99, and the extreme-value threshold fraction among 5, 10 and 15 percent. Each is chosen by
the lowest summed average FZ0 loss, over the two selection series (the equity market and the 10-year bond) and the two tail levels, of the model that
uses it, scored on every \mkSpanStride{}th forecast day of the span \mkSpanStart{} to \mkSpanEnd{} (every second one of those for the extreme-value model). \Cref{tab:market-selection} lists every candidate
with its score. The euro series starts in 2000 and cannot carry a window inside the span, so the foreign-exchange book is not part of the
selection. The selected values are an estimation window of \mkSelEstimationWindowDays{} days, an EWMA decay of \mkSelEwmaLambda{}, an
age-weighting decay of \mkSelAgeWeightLambda{} and a threshold fraction of \mkSelEvtThresholdFraction{}. The numbers of CAViaR and GAS starts, the
refit intervals (\mkCaviarRefit{} forecast days for CAViaR, \mkGasRefit{} for the GAS-FZ models, every day for the GARCH family, the filtered
simulation, the extreme-value model and the ES step of CAViaR) and the seed \mkSeed{} are fixed by the compute limit and the determinism
requirement and are not tuned.

\begin{table}[htbp]
\centering
\footnotesize
\setlength{\tabcolsep}{3.0pt}
\caption{The hyperparameters fixed once on 1995 to 2003, before the earliest test window: for each candidate value the sum of the average FZ0 loss of the model that uses it over the two selection series and the two tail levels, on every fifth forecast day of the span, and the value chosen. The last two rows are the quantile network.}
\label{tab:market-selection}
\begin{tabular}{llp{6.0cm}r}
\toprule
Hyperparameter & Model scored & Candidate: summed average FZ0 (lower is better) & Selected \\
\midrule
\texttt{age\_weight\_lambda} & Age-weighted HS & 0.98: -14.36, 0.99: -14.56 & 0.99 \\
\texttt{estimation\_window\_days} & GARCH-t & 500: -15.10, 750: -15.17, 1000: -15.18 & 1000 \\
\texttt{evt\_threshold\_fraction} & Cond. EVT & 0.05: -16.12, 0.1: -16.11, 0.15: -16.11 & 0.05 \\
\texttt{ewma\_lambda} & EWMA & 0.94: -14.80, 0.97: -14.81, 0.99: -14.53 & 0.97 \\
\texttt{hidden\_units} & Quantile network & 4: -7.25, 8: 0.63, 16: 19.35 & 4 \\
\texttt{weight\_decay} & Quantile network & 0.001: 0.63, 0.01: -6.95, 0.1: -14.93 & 0.1 \\
\bottomrule
\end{tabular}
\end{table}


\Cref{tab:market-selection} is read in two ways. The scores of the estimation window differ by \mkSelWinGap{} between 500 and 1,000 days, so the choice of
the longest window rests on a small margin; the score of the EWMA decay is flat between 0.94 and 0.97 and worse at 0.99; the scores of the threshold
fraction of the extreme-value model differ in the third decimal, so that its selected value is a marginal win. What the table does not show
is any sensitivity of the test-window results to these choices, because the hyperparameters are not revisited after the test windows.

The cost of an estimator is set by the length $N$ of the window, the number $T$ of forecast days, the refit interval $R$ and the work
of one fit. Historical simulation and the age-weighted variant sort the window each day, $O(TN\log N)$ in all; the delta-normal and
Cornish--Fisher estimators compute moments of the window, $O(TN)$; the EWMA is $O(N+T)$ because the recursion updates in constant time. A
GARCH-family fit evaluates the likelihood recursion, which costs $O(N)$ per evaluation, $I$ times in the optimiser, and runs $T/R$ times: $O\big((T/R)\,I\,N\big)$, with
$R=1$ in the case. Filtered historical simulation and the extreme-value model add a sort of the residuals per day. A CAViaR or a GAS-FZ fit
costs $O(SIN)$ for $S$ starts and $I$ evaluations per start, every $R=\mkCaviarRefit{}$ or $\mkGasRefit{}$ days, and the objective is non-smooth, so
$I$ is large. The scoring is $O(T)$ per model and level.

The engine implements every estimator of this subsection in the \texttt{riskvalidation} package, release 1.1.0:
\mkfn{riskvalidation.market.historical} (the sample estimator, age-weighted simulation, filtered and volatility-updated simulation),
\mkfn{riskvalidation.market.parametric} (delta-normal, Cornish--Fisher, EWMA), \mkfn{riskvalidation.market.garch} and
\mkfn{riskvalidation.market.evt} (the GARCH family through the optional \texttt{arch} package, the generalised Pareto tail),
\mkfn{riskvalidation.market.caviar}, \mkfn{riskvalidation.market.gas_fz}, \mkfn{riskvalidation.market.scoring} (FZ0 and tick loss) and
\mkfn{riskvalidation.market.conventions} (the sign convention). Each function receives no data of the day it forecasts. The pipeline of
case C11 only slices the window, schedules the refits and collects the numbers, in \mkfn{data-pipeline/pipeline/cases/c11/forecast.py}
(\mkfn{run_model}) and \mkfn{hyper.py}. The acceptance criterion of the product is that each estimator runs over the full rolling schedule
with no look-ahead, that the forecast of the Student-$t$ GARCH agrees, at fixed parameters, with the filter of the \texttt{arch} package within a relative $10^{-9}$ and its closed-form expected shortfall with the quadrature of the definition within $10^{-4}$ (the fit itself is the \texttt{arch} package, so no second likelihood is compared; the test is \mkfn{test_garch_forecast_student_t_closed_form}), and that FZ0 losses are computed for every
estimator that forecasts an ES. The estimators are compared with one another on the same days by the tests of \cref{sec:m19}.

% ----------------------------------------------------------------------------------------------------------------------------------
\subsection{Back-tests of value at risk and expected shortfall}
\label{sec:m19}

\subsubsection{What a back-test asks}
A back-test confronts a sequence of forecasts with the outcomes. \citet{market:noldeziegel2017} separate two kinds. A traditional
back-test takes the hypothesis that the model is correct and asks whether the realised sequence is compatible with it: the coverage,
independence and tail tests below. A comparative back-test takes two forecasters and asks whether the internal model is at least as good as
a reference, through a consistent scoring function: the Diebold--Mariano and model confidence tests below. A validator wants both, because
a rejected model is not a better one and a passed model has not been shown to beat another. Let $\mathcal I_t=\indic\{\ell_t>\mathrm{VaR}_t\}$ be the
exceedance indicator of day $t$, where $\ell_t=-y_t$ is the loss and VaR is stated as a positive loss number. A correct model has
conditional coverage, which means that
\begin{equation}
\mathcal I_t\mid\mathcal F_{t-1}\;\sim\;\mathrm{Bernoulli}(\alpha)\quad\text{independently},\qquad x=\sum_{t=1}^T\mathcal I_t,
\label{eq:mk-exceed}
\end{equation}
and every coverage test is a test of this statement or of one of its consequences: the right frequency, no dependence between
exceedances, no predictability of the hits. A test reports a statistic, a $p$-value and a light. The lights are screening levels of the
validation, not regulatory thresholds: red when $p<\mkRedAlpha{}$, amber when $p<\mkAmberAlpha{}$, green otherwise. A test that cannot run, for
example the independence test with no exceedance, is reported as not evaluated with its reason and never as a number; the convention
$0^0=1$ would return the value zero for a test that did not run.

\subsubsection{Coverage, independence and duration}
The proportion-of-failures test of \citet{market:kupiec1995} compares the observed frequency $\widehat\pi=x/T$ with $\alpha$ by the likelihood ratio of
two Bernoulli models, which is asymptotically chi-square with one degree of freedom,
\begin{equation}
LR_{POF}=-2\ln\frac{(1-\alpha)^{T-x}\alpha^{x}}{(1-\widehat\pi)^{T-x}\widehat\pi^{x}}\;\overset{d}{\longrightarrow}\;\chi^2_1.
\label{eq:mk-pof}
\end{equation}
The implementation also reports the exact binomial $p$-value, the total probability of the outcomes at most as likely as $x$ under
$\mathrm{Bin}(T,\alpha)$. With $T=250$ and $\alpha=0.01$ the expected count is $2.5$, the power is low, and the asymptotic test is not of nominal size:
at the 5 percent level, where the light turns red, it rejects when $x=0$ or $x\ge7$, 9.5 percent of the time under a correct model, and at the 10 percent level, where
it turns amber, when $x=0$ or $x\ge6$, 12.2 percent of the time, so the exact value is the safe reading at that window
while the light follows the asymptotic $p$-value. In the windows of case C11 the expected counts are between \mkExpMixed{} and \mkExpCalm{} at 99 percent
(\cref{tab:market-exceed}) and the exact value changes few readings.

The independence test of \citet{market:christoffersen1998} models the exceedance sequence as a first-order Markov chain. With $n_{ij}$ the number of
transitions from state $i$ to state $j$, $\pi_{01}=n_{01}/(n_{00}+n_{01})$, $\pi_{11}=n_{11}/(n_{10}+n_{11})$ and
$\pi=(n_{01}+n_{11})/\sum n_{ij}$,
\begin{equation}
LR_{ind}=-2\ln\frac{(1-\pi)^{n_{00}+n_{10}}\pi^{n_{01}+n_{11}}}{(1-\pi_{01})^{n_{00}}\pi_{01}^{n_{01}}(1-\pi_{11})^{n_{10}}\pi_{11}^{n_{11}}}\;\overset{d}{\longrightarrow}\;\chi^2_1,
\qquad LR_{cc}=LR_{POF}+LR_{ind}\;\overset{d}{\longrightarrow}\;\chi^2_2.
\label{eq:mk-ind}
\end{equation}
The conditional coverage statistic $LR_{cc}$ tests the frequency and the independence together. It is not evaluated when there is no
exceedance or every day is an exceedance, because nothing is left to test for clustering. The duration test of
\citet{market:christoffersenpelletier2004} models the number of days between exceedances with the Weibull hazard $\lambda(d)=a^bb\,d^{b-1}$; under
correct conditional coverage the durations are memoryless and $b=1$. With $\mathcal C$ the set of the $d_c$ completed spells $D_i$ between consecutive
exceedances and $\mathcal K$ the set of the censored spells, the first and the last,
\begin{equation}
LR=2\left[\ell(\widehat a,\widehat b)-\ell(\widetilde a,1)\right],\qquad
\ell(b)=d_c\ln\frac{d_c}{\sum_{i\in\mathcal C\cup\mathcal K}D_i^{\,b}}+d_c\ln b+(b-1)\sum_{i\in\mathcal C}\ln D_i-d_c,
\label{eq:mk-dur}
\end{equation}
where $a^b=d_c/\sum_{i\in\mathcal C\cup\mathcal K}D_i^b$ is profiled out. A censored spell enters through its survival term $\exp\!\big(-(aD)^b\big)$ only, so it adds to the
denominator and not to the sums of the logarithms or to the count $d_c$ (the code is \mkfn{duration_test}). The $p$-value is by Monte Carlo: $n_{mc}=\mkMcDraws{}$ Bernoulli sequences of the same length give
$p=\big(1+\#\{LR_{sim}\ge LR\}\big)/(n_{mc}+1)$, and a simulated sequence with fewer than two exceedances has the statistic zero. Fewer than two
exceedances in the data make the test not evaluated. The dynamic quantile test of \citet{market:engleManganelli2004} regresses the hit
$\mathrm{Hit}_t=\mathcal I_t-\alpha$ on a constant, its four lags and the VaR itself, $X_t=(1,\mathrm{Hit}_{t-1},\dots,\mathrm{Hit}_{t-4},\mathrm{VaR}_t)$,
\begin{equation}
DQ=\frac{\mathrm{Hit}'X\,(X'X)^{+}X'\,\mathrm{Hit}}{\alpha(1-\alpha)},
\label{eq:mk-dq}
\end{equation}
asymptotically chi-square with the rank of $X$ as degrees of freedom, and invariant to an affine change of the VaR regressor, which leaves the column space of $X$ unchanged. The asymptotic law is poor with sparse hits,
so the light follows a Monte Carlo $p$-value of the same form as the duration test. \Cref{tab:market-tests} lists the null hypothesis and the
alternative of each test as the product records them.

\begin{table}[htbp]
\centering
\scriptsize
\setlength{\tabcolsep}{2.5pt}
\caption{The six tests the case C11 battery runs on every model at every tail level, with the null hypothesis and the alternative as the committed test records state them and the p-value below which the light turns amber and red; the lights are internal screening levels of the validation, not regulatory thresholds.}
\label{tab:market-tests}
\begin{tabular}{p{2.9cm}p{4.4cm}p{4.4cm}rr}
\toprule
Test & Null hypothesis & Alternative & Amber below & Red below \\
\midrule
\relax Kupiec proportion of failures & The exceedance probability equals alpha & the exceedance probability differs from alpha & 0.1 & 0.05 \\
\relax Christoffersen independence & The exceedances are independent over time & an exceedance is more (or less) likely after an exceedance & 0.1 & 0.05 \\
\relax Christoffersen conditional coverage & The exceedance probability equals alpha and the exceedances are independent & the coverage is wrong or the exceedances cluster & 0.1 & 0.05 \\
\relax Duration test & The no-hit durations are memoryless (Weibull b = 1) & the durations have a Weibull hazard with b != 1 (clustering or too regular exceedances) & 0.1 & 0.05 \\
\relax Dynamic quantile test & The hits are unpredictable from their past and from the VaR (conditional coverage holds) & the hits are predictable: the quantile forecasts are conditionally biased & 0.1 & 0.05 \\
\relax Acerbi-Szekely $Z_2$ & The ES forecasts are correct (the statistic has the simulated null distribution) & risk is underestimated (the statistic is negative) & 0.05 & 0.0001 \\
\bottomrule
\end{tabular}
\end{table}


\subsubsection{The traffic light}
The Basel zones are the cumulative binomial law of the exception count at true coverage. For $x$ exceptions in $n=250$ days at the 99 percent level,
\begin{equation}
P(X\le x)=\sum_{i=0}^{x}\binom{n}{i}\,0.01^{\,i}\,0.99^{\,n-i},\qquad X\sim\mathrm{Bin}(250,0.01),
\label{eq:mk-binom}
\end{equation}
and the zone design of the Committee places the amber zone where this probability first reaches 95 percent and the red zone where it reaches
99.99 percent \citep{market:bcbs1996}: green for 0 to 4 exceptions, amber for 5 to 9, red for 10 or more, with the same rule giving the zones
of another window length. \Cref{tab:market-traffic} reproduces the zones and multipliers of MAR32.9 \citep{market:bcbsmar32} and the cumulative probabilities of MAR99.17 (Table 2) \citep{market:bcbsmar99} as the models document
records them, with the probabilities of \mkTlFourCum{} percent at four exceptions, \mkTlFiveCum{} at five, \mkTlNineCum{} at nine and \mkTlTenCum{} at ten,
the backtesting-dependent multipliers from \mkTlFourMult{} to \mkTlTenMult{}, and the factors of the United States
rule from \mkTlFourUs{} to \mkTlTenUs{} \citep{market:cfr217204}. The committed probabilities agree with the binomial law of
equation~\eqref{eq:mk-binom} to two decimals (\cref{fig:market-binomial}). A window shorter than 250 forecast days carries the light as not
evaluated with its reason, because the light is defined on 250 observations.

\begin{table}[htbp]
\centering
\footnotesize
\setlength{\tabcolsep}{3.0pt}
\caption{The traffic-light table of the 99 percent value-at-risk back-test over 250 observations as the committed models document records it: the zone, the backtesting-dependent multiplier of the Basel framework, the multiplication factor of the United States rule and the cumulative binomial probability of the exception count at true coverage.}
\label{tab:market-traffic}
\begin{tabular}{rlrrr}
\toprule
Exceptions in 250 days & Zone & FRTB multiplier & US multiplication factor & Cumulative probability (\%) \\
\midrule
0 & green & 1.50 & 3.00 & 8.11 \\
1 & green & 1.50 & 3.00 & 28.58 \\
2 & green & 1.50 & 3.00 & 54.32 \\
3 & green & 1.50 & 3.00 & 75.81 \\
4 & green & 1.50 & 3.00 & 89.22 \\
5 & amber & 1.70 & 3.40 & 95.88 \\
6 & amber & 1.76 & 3.50 & 98.63 \\
7 & amber & 1.83 & 3.65 & 99.60 \\
8 & amber & 1.88 & 3.75 & 99.89 \\
9 & amber & 1.92 & 3.85 & 99.97 \\
10 & red & 2.00 & 4.00 & 99.99 \\
\bottomrule
\end{tabular}
\end{table}


\begin{figure}[htbp]
\centering
\includegraphics[width=0.95\textwidth]{market-binomial.pdf}
\caption{Binomial law of the number of exceptions in 250 days at 99 percent coverage (bars, hatched by the Basel zone) and its cumulative
probability (line, right axis, with the 95 and 99.99 percent levels that start the amber and the red zone). The probabilities are computed from the
binomial distribution and checked against the committed table of the models document.}
\label{fig:market-binomial}
\end{figure}

\subsubsection{Expected shortfall tests}
The Acerbi--Szekely tests compare the loss beyond the VaR with the forecast ES. With $X_t=-\ell_t$ the profit and loss, $N_T=\sum_t\mathcal I_t$ the number of exceedances and
$\ES_t>0$ the forecast as a positive loss number,
\begin{equation}
Z_1=\frac1{N_T}\sum_{t=1}^{T}\frac{X_t\,\mathcal I_t}{\ES_t}+1,\qquad
Z_2=\sum_{t=1}^{T}\frac{X_t\,\mathcal I_t}{T\alpha\,\ES_t}+1,
\label{eq:mk-az}
\end{equation}
both of expectation zero under the null and negative when risk is underestimated \citep{market:acerbiszekely2014}. The critical values depend on the
model, so they are simulated: with $n_{sim}=\mkZDraws{}$ loss paths $Z_{sim}$ drawn from a null law, $p=(1+\#\{Z_{sim}\le Z\})/(n_{sim}+1)$, with the light taken against
$\mkAzAmber{}$ for amber and $\mkAzRed{}$ for red, a level that needs at least 10,000 draws to resolve. A historical, CAViaR or GAS model has no predictive
distribution from which to draw, so case C11 uses one null for the whole ladder. On each day the loss exceeds the VaR with probability $\alpha$,
and by an exponential excess whose mean is the gap between the ES and the VaR,
\begin{equation}
\ell^{sim}_t=\mathrm{VaR}_t+\mathcal E_t\,(\ES_t-\mathrm{VaR}_t)\ \text{ with probability }\alpha,\qquad\ell^{sim}_t=0\ \text{ otherwise},\qquad\mathcal E_t\sim\mathrm{Exp}(1).
\label{eq:mk-null}
\end{equation}
The simulated paths reproduce each model's own VaR and ES, so the test asks only whether the ES level is right; the exponential shape (the generalised Pareto with shape
zero) is an assumption, and a heavier true tail makes the critical values too narrow. The reference point of $-\mkZRefYellow{}$ (yellow) and $-\mkZRefRed{}$ (red)
for $T=\mkZRefN{}$ and $\alpha=\mkZRefAlpha{}$ percent is a published suggestion that the product verifies through the discussion of \citet{market:moldenhauerpitera2017}
only; it is carried with the result and is never a threshold. The Acerbi--Szekely statistic is not evaluated when an ES forecast is undefined, as with
a conditional extreme-value forecast whose shape reaches one, and the $Z_1$ statistic needs at least one exceedance.

Two further tests are in the engine and are not run on C11, because they need the probability integral transform of each day or a volatility path
that not every estimator has. The test of \citet{market:duescanciano2017} works on the transform $u_t=F_t(Y_t)$ through the cumulative violation
$H_t=\frac1\alpha(\alpha-u_t)\indic\{u_t\le\alpha\}$, which under the null has mean $\alpha/2$ and variance $\alpha(1/3-\alpha/4)$ when $u$ is uniform
\citep{market:huehurlinlu2024},
\begin{equation}
U_{ES}=\frac{\sqrt T\left(\bar H-\alpha/2\right)}{\sqrt{\alpha\left(1/3-\alpha/4\right)}}\;\overset{d}{\longrightarrow}\;N(0,1),\qquad
C_{ES}(m)=T\sum_{j=1}^{m}\widehat\rho_j^{\,2}\;\overset{d}{\longrightarrow}\;\chi^2_m,
\label{eq:mk-de}
\end{equation}
where $\widehat\rho_j$ is the sample autocorrelation of $H_t$ at lag $j$. A large positive $U_{ES}$ means underestimated risk, so the $p$-value is the
one-sided upper tail. The residual test of \citet{market:mcneilfrey2000} takes the exceedance days, on which
$r_t=(\ell_t-\ES_t)/\sigma_t$ should be independent with mean zero, and tests the mean by a one-sided bootstrap of the centred residuals. Other
approaches exist: a spectral-measure back-test \citep{market:costanzinocurran2015} and multinomial back-tests of the VaR at several levels as an implicit
test of the ES \citep{market:kratzlokmcneil2018}. The plain asymptotic versions ignore parameter estimation uncertainty. The size and power of every test
of the section, measured on data with a known truth, are in \cref{tab:market-sizepower}.

\subsubsection{Comparison, eligibility and the champion}
The comparative evidence is the average FZ0 loss of equation~\eqref{eq:mk-fz0} over the forecast days, a lower value meaning a better forecast of the pair.
The difference of the losses of two models, $d_t=L_t^{A}-L_t^{B}$, is tested by the statistic of \citet{market:dieboldmariano1995} with a
long-run variance estimated with the Bartlett kernel and $h$ lags (heteroskedasticity and autocorrelation consistent for $h>0$),
\begin{equation}
DM=\frac{\bar d}{\sqrt{\widehat{LRV}/T}},\qquad
\widehat{LRV}=g_0+2\sum_{j=1}^{h}\Big(1-\frac j{h+1}\Big)g_j,\qquad g_j=\frac1T\sum_t(d_t-\bar d)(d_{t-j}-\bar d),
\label{eq:mk-dm}
\end{equation}
with $p=2(1-\Phi(|DM|))$; a negative mean difference means that model $A$ has the lower loss. The forecasts are one day ahead, so
the case uses $h=\mkHacLags{}$ lags, and with $h=0$ the long-run variance is the sample variance of $d_t$. For the VaR alone the loss is the tick loss, and for the pair it is FZ0. The model confidence set of
\citet{market:hansenlundenason2011} eliminates models sequentially by an equivalence test with a stationary bootstrap and returns the set that contains the
best model with the stated confidence; the case uses the range statistic $T_R$, a confidence of \mkMcsConf{} percent, \mkMcsBoot{} bootstrap
resamples and blocks of \mkMcsBlock{} days, on the models that have an ES on every day. The set is never empty, and it is large when the score cannot separate
close models. It is a comparative back-test in the sense of \citet{market:noldeziegel2017}, not a proof that the best member is right.

The two kinds of evidence are combined by a rule that every model and window obey. A model is eligible when none of the Kupiec, the conditional coverage
and the Acerbi--Szekely tests is red at either tail level, and the champion at level $\alpha$ is the eligible model with an ES on every day and the lowest average
FZ0 loss,
\begin{equation}
\mathrm{champion}(\alpha)=\arg\min_{m\,\in\,\mathcal M_{\mathrm{elig}}}\ \frac1T\sum_{t=1}^{T}L_{FZ0}\big(y_t,\widehat v^{\,m}_t,\widehat e^{\,m}_t;\alpha\big),
\label{eq:mk-champion}
\end{equation}
where $\mathcal M_{\mathrm{elig}}$ is the set of eligible models that forecast an ES on every day; when the set is empty the level has no champion, and the
reason is recorded. Each other model is compared with the champion by the statistic of equation~\eqref{eq:mk-dm}. \Cref{alg:mk-rolling} states the whole
procedure of one window.

\begin{algorithm}[htbp]
\caption{Rolling evaluation of one window of case C11: forecasts, back-tests, eligibility and champion.}
\label{alg:mk-rolling}
\begin{algorithmic}[1]
\Require returns $y$, forecast days $t=t_0,\dots,t_0+T-1$, models $\mathcal M$, levels $\alpha\in\{0.01,0.025\}$, hyperparameters fixed on the selection span, seed
\For{each model $m\in\mathcal M$ and each day $t$}
\State $W\gets y_{t-N},\dots,y_{t-1}$ (the window of the model; nothing of day $t$ or later)
\If{a refit of $m$ is due at $t$}
\State refit $m$ on $W$; if the fit fails, keep the previous fit and count the failure
\EndIf
\State run the filter of $m$ over $W$; store $(\widehat v^m_t,\widehat e^m_t)$ at each $\alpha$ in the return convention; hold the previous forecast if none is usable and count it
\EndFor
\State convert to positive loss numbers; round every array as it is stored so that each score recomputes from the stored arrays
\For{each model $m$ and each level $\alpha$}
\State $\mathcal I_t\gets\indic\{\ell_t>\mathrm{VaR}^m_t\}$; run the Kupiec, independence, conditional coverage, duration, dynamic quantile and Acerbi--Szekely tests of this subsection
\State Basel zone over every rolling window of 250 days (not evaluated below 250 forecast days); average FZ0 and tick loss
\EndFor
\State $\mathrm{eligible}(m)\gets$ no red Kupiec, conditional coverage or Acerbi--Szekely light at either level
\For{each level $\alpha$}
\State champion $\gets\arg\min$ average FZ0 over eligible models with an ES on every day (none if the set is empty)
\State Diebold--Mariano statistic of every model against the champion; model confidence set over the models with an ES on every day
\EndFor
\State \Return the arrays, the test rows, the eligibility flags, the champion and the comparison
\end{algorithmic}
\end{algorithm}

\subsubsection{The desk tests of the framework and the desk generator}
The two desk tests are the desk back-test and the profit-and-loss attribution test. The desk back-test counts the days of the last 250 on which the hypothetical loss
exceeds the desk VaR, at 99 and at 97.5 percent, and flags the desk when the count exceeds the limit of \mkDeskLimitNN{} or \mkDeskLimitNF{} respectively; a missing
P\&L or VaR counts as an exception for the bank-wide test (MAR32.5) and for a desk (MAR32.18), and the product implements both through \mkfn{desk_exceptions} with the missing values counted;
the flag \mkfn{missing_as_exception} (true by default) that turns the count off is a sensitivity, not the rule. The attribution test compares,
over the same 250 days, the hypothetical P\&L $R_{HPL}$ and the risk-theoretical P\&L $R_{RTPL}$ through
\begin{align}
r_S&=\frac{\mathrm{cov}(\mathrm{rk}\,R_{HPL},\mathrm{rk}\,R_{RTPL})}{\sigma_{\mathrm{rk}R_{HPL}}\,\sigma_{\mathrm{rk}R_{RTPL}}},\\
KS&=\sup_x\big|F_{RTPL}(x)-F_{HPL}(x)\big|,\qquad F(x)=0.004\,\#\{\text{observations}\le x\},
\label{eq:mk-pla}
\end{align}
the Spearman correlation as the Pearson correlation of the average ranks and the Kolmogorov--Smirnov distance of the two empirical distributions, with the zone: green when
$r_S>\mkThrSpGreen{}$ and $KS<\mkThrKsGreen{}$, red when $r_S<\mkThrSpRed{}$ or $KS>\mkThrKsRed{}$, amber otherwise, the boundaries belonging to amber \citep{market:bcbsmar32}. A red desk is
ineligible for the internal models approach, and an amber desk stays on it with a capital surcharge. The distance is computed as an exact integer count of steps of the empirical
distribution functions, $KS=\text{gap}/250$, and the zone on the integer ranks, so that a value on a boundary is never moved by a floating-point quotient: a gap of 30 steps is exactly $0.12$
and is amber.

To measure what these tests catch, case C13 simulates a desk of which the truth is known. A two-factor spread desk is long a factor with the loading $l_1=\mkGLoadA{}$ and short a
factor with the loading $l_2=\mkGLoadB{}$ per unit standard deviation; the factors have the correlation $\rho=\mkGRho{}$. The day scale is a log-AR(1) known at the start of the day, so that a VaR
built on it is an ex-ante forecast, with persistence $\phi=\mkGPhi{}$ and stationary log standard deviation $s=\mkGVolSd{}$, and the constant $\mu_\sigma=-s^2$ gives $\E\sigma^2=1$:
\begin{equation}
\sigma_t=\exp\!\left(\mu_\sigma+g_{t-1}\right),\qquad g_t=\phi\,g_{t-1}+s\sqrt{1-\phi^2}\,\epsilon_t,\qquad
\mathrm{HPL}_t=\sigma_t\,l'x_t+\kappa\sqrt{l'Rl}\;u_t,
\label{eq:mk-gen}
\end{equation}
with $x_t$ the correlated factor shocks of correlation matrix $R$, $\epsilon_t$ and $u_t$ independent standard normals, $x_t$ drawn from the first two columns of $w_t$, a vector of three independent standard normals (the third column is the omitted factor of one deficiency), $b_t$ a further independent standard normal, the basis shock added to the hypothetical P\&L in the deficiency that leaves a basis out, and $\kappa=\mkGResShare{}$ the share of the factor P\&L
standard deviation that is pricing noise. The risk-theoretical P\&L is the risk model's revaluation of the same moves, and the desk VaR is $\Phi^{-1}(1-\alpha)$ times the model's
conditional standard deviation of the factor P\&L. A deficiency changes only the risk model or the date of its P\&L (\cref{tab:market-truth}): the model scales the day's shocks with
the scale of 20 days before; it uses a correlation 0.3 too high; the hypothetical P\&L carries a basis or a third factor that the model leaves out, sized as a share of its
standard deviation; or the series is dated one day off. The population Pearson correlation of the two series follows from the parameters. For a lag of $k$ days the cross moment of the two
scales is $\E[\sigma_t\sigma_{t-k}]=\exp\!\left(2\mu_\sigma+s^2(1+\phi^k)\right)$ because $g$ is Gaussian and stationary, which equals
\begin{equation}
\E\left[\sigma_t\sigma_{t-k}\right]=\exp\!\left(-s^2\left(1-\phi^{k}\right)\right),
\label{eq:mk-pearson}
\end{equation}
so that a stale scale lowers the correlation by this factor and leaves the marginal distribution of the risk-theoretical P\&L equal to that of the unlagged series, since $\sigma_{t-k}$ and $\sigma_t$
have the same distribution. The closed form is what lets the generator check its simulation. \Cref{alg:mk-desk} states the measurement of one variant.

\begin{table}[htbp]
\centering
\footnotesize
\setlength{\tabcolsep}{3.0pt}
\caption{The six variants of the desk generator of case C13: the planted parameter, the statistic the deficiency is meant to break, the population Pearson correlation of the two P\&L series in closed form from the generator parameters, and the mean Spearman correlation measured over 1,000 replications.}
\label{tab:market-truth}
\begin{tabular}{lp{4.0cm}lrr}
\toprule
Variant & Planted parameter & Statistic meant to fail & Population Pearson & Mean Spearman \\
\midrule
None (size) & -- & none & 0.995 & 0.992 \\
Missing basis & basis sd share 0.65 & spearman & 0.756 & 0.705 \\
Stale volatility & vol lag days 20 & ks & 0.955 & 0.979 \\
Wrong correlation & correlation error 0.3 & either & 0.990 & 0.986 \\
Omitted factor & omitted loading share 0.6 & spearman & 0.796 & 0.772 \\
Time misalignment & lag days 1 & spearman & 0.000 & -0.004 \\
\bottomrule
\end{tabular}
\end{table}


\begin{algorithm}[htbp]
\caption{Measurement of one variant of the desk case C13: replications through the engine.}
\label{alg:mk-desk}
\begin{algorithmic}[1]
\Require variant, case seed, number of replications $n_{rep}$, window $250$ days
\For{$i=0,\dots,n_{rep}-1$}
\State seed $\gets$ seed base of the variant $+\,i$; draw $300+250$ days of $\epsilon,w,u,b$ in a fixed order; keep the last 250 days
\State build $\mathrm{HPL}$, $\mathrm{RTPL}$ and the desk VaR at 99 and 97.5 percent with the planted deficiency, if any
\State $r_S\gets$ Spearman$(\mathrm{HPL},\mathrm{RTPL})$; $KS\gets$ Kolmogorov--Smirnov$(\mathrm{HPL},\mathrm{RTPL})$; joint zone; zone of each statistic with the other held at its best value
\State exceptions of the hypothetical loss against the desk VaR at each level
\EndFor
\State \Return mean, standard deviation and quantiles of $r_S$ and $KS$, the zone counts, the exception summaries and the Wilson interval of each flagged share
\end{algorithmic}
\end{algorithm}

\subsubsection{Parameters, complexity and implementation}
The parameters of the battery are the tail levels $\alpha=0.01$ and $\alpha=0.025$, the 5 percent significance of the comparison, the screening levels of \cref{tab:market-tests}, \mkZDraws{} simulated paths
for the Acerbi--Szekely null and \mkMcDraws{} for the duration and the dynamic quantile tests, the lag count 4 of the dynamic quantile regression, and the bootstrap of the confidence set of
\cref{sec:m19}. They are conventions of the validation, stated with every result, and none is tuned on a window. The statistics $LR_{POF}$ and $LR_{ind}$ cost $O(T)$; the duration test
costs $O(T)$ per likelihood evaluation plus $O(n_{mc}T)$ for the simulation; the Acerbi--Szekely test costs $O(n_{sim}T)$; the model confidence set resamples the $T\times M$ loss matrix $B$ times at each elimination step, so that
its cost grows at least as $O(B\,T\,M)$ for $B$ bootstrap resamples and $M$ models; the Diebold--Mariano statistic costs $O(Th)$. In case C13 the cost is the replication count times $O(250)$ per statistic, with the
Spearman statistic dominated by a sort.

The tests are implemented in \mkfn{riskvalidation.validation.market} (\mkfn{kupiec_pof}, \mkfn{christoffersen_independence},
\mkfn{christoffersen_conditional_coverage}, \mkfn{duration_test}, \mkfn{dynamic_quantile_test}, \mkfn{basel_traffic_light}, \mkfn{acerbi_szekely},
\mkfn{du_escanciano_unconditional}, \mkfn{du_escanciano_conditional}, \mkfn{mcneil_frey_residual_test}, \mkfn{pla_test}, \mkfn{desk_backtest}), with the scoring
and the model confidence set in \mkfn{riskvalidation.market.scoring} and the desk statistics and zones in \mkfn{riskvalidation.market.frtb}, all in release 1.1.0 of the
\texttt{riskvalidation} package. Each returns one result record with its statistic, $p$-value, light and reference. The pipeline modules are
\mkfn{data-pipeline/pipeline/cases/c11/documents.py} (the tests, the null of equation~\eqref{eq:mk-null}, the eligibility and the comparison) and
\mkfn{data-pipeline/pipeline/cases/c13/generator.py} and \mkfn{measure.py} (the desk generator and the measurement). The acceptance criterion of the product is that the traffic-light table of MAR99 (Table 2)
is reproduced to two decimals (\cref{fig:market-binomial}), that the Kupiec and Christoffersen statistics agree with their closed forms, and that the Acerbi--Szekely critical values are simulated per
model; the size and power of the tests on data with a known truth are the evidence that each test behaves as its theory says.

\begin{table}[htbp]
\centering
\scriptsize
\setlength{\tabcolsep}{2.5pt}
\caption{Size and power of the nine back-tests measured by the validator of case C22: the rejection rate at the 5 percent level when the null holds at its boundary and against one stated alternative, over two thousand repetitions each, with the Wilson 95 percent interval in brackets; the size of a test that rejects at its red or Monte Carlo level is read at that level.}
\label{tab:market-sizepower}
\begin{tabular}{p{2.9cm}p{6.6cm}rr}
\toprule
Test & Sample, level and alternative & Size (\%) & Power (\%) \\
\midrule
\relax McNeil-Frey residual test & 2,500 days at 5\%, 125 exceedances expected; the ES forecast is 95\% of the true 5\% expected shortfall of a normal loss (mean residual 0.103) & 4.5 [3.6, 5.4] & 93.2 [92.0, 94.2] \\
\relax Kupiec proportion of failures & 1,500 days at 5\%; the VaR is exceeded 1.3 times too often (a hit rate of 6.5\%, not 5\%) & 5.1 [4.2, 6.2] & 69.3 [67.2, 71.3] \\
\relax Christoffersen independence & 1,500 days at 5\%; a hit is followed by a hit with probability 15\%, not 5\% (clustered exceedances at the same hit rate) & 6.6 [5.6, 7.8] & 87.5 [86.0, 88.9] \\
\relax Christoffersen conditional coverage & 1,500 days at 5\%; a hit is followed by a hit with probability 15\%, not 5\% (clustered exceedances at the same hit rate) & 6.7 [5.6, 7.8] & 81.2 [79.4, 82.9] \\
\relax Duration test & 400 days at 5\%; a hit is followed by a hit with probability 40\%, not 5\% (clustered exceedances at the same hit rate) & 5.4 [4.5, 6.5] & 63.8 [61.7, 65.9] \\
\relax Acerbi-Szekely $Z_2$ & 250 days at 2.5\%; the realised losses are 1.25 times as dispersed as the VaR and ES forecasts say & 4.8 [3.9, 5.8] & 89.4 [88.0, 90.7] \\
\relax Du-Escanciano, unconditional & 2,000 days at 10\%; the lower tail has too much mass (u = r\textasciicircum{}1.05: 11.2\% of the PITs below 0.10, not 10\%) & 4.7 [3.9, 5.7] & 56.5 [54.3, 58.7] \\
\relax Du-Escanciano, conditional & 2,000 days at 10\%; the tail PITs (below 0.10) cluster: one follows another with probability 20\%, not 10\% & 5.5 [4.6, 6.6] & 73.2 [71.2, 75.1] \\
\relax Dynamic quantile test & 1,500 days at 5\%; a hit is followed by a hit with probability 15\%, not 5\% (clustered exceedances at the same hit rate) & 4.8 [3.9, 5.8] & 73.7 [71.7, 75.6] \\
\relax Dynamic quantile test & 250 days at 1\%; a hit is followed by a hit with probability 10\%, not 1\% (clustered exceedances at the same hit rate) & 5.1 [4.2, 6.1] & 16.0 [14.4, 17.6] \\
\relax Dynamic quantile test & 250 days at 2.5\%; a hit is followed by a hit with probability 15\%, not 2.5\% (clustered exceedances at the same hit rate) & 5.9 [5.0, 7.1] & 25.7 [23.8, 27.7] \\
\relax Dynamic quantile test & 750 days at 1\%; a hit is followed by a hit with probability 10\%, not 1\% (clustered exceedances at the same hit rate) & 4.2 [3.4, 5.1] & 27.9 [26.0, 29.9] \\
\bottomrule
\end{tabular}
\end{table}


% ----------------------------------------------------------------------------------------------------------------------------------
\subsection{Quantile network trained on the FZ0 loss}
\label{sec:m30}

\subsubsection{Why a network on the joint loss}
The joint elicitability of the pair $(v,e)$ \citep{market:fisslerziegel2016} has a practical consequence that the econometric estimators of \cref{sec:m18} use only partly: the pair can be estimated by minimising
the average FZ0 loss of equation~\eqref{eq:mk-fz0} over a flexible family of functions of the available information. The GAS-FZ models do so with a recursion of one or two
factors; a neural network does so with a function of the recent returns and volatilities that is not constrained to a recursion. A line of recent work trains networks for the joint VaR
and ES; the large model of \citet{market:zhang2026resga} learns the pair on cross-sections of firm characteristics, which is a different problem from the forecast of one portfolio series. The
model of C11 is the compact single-series network of that family: one hidden layer, trained on the FZ0 loss of the series' own window, and not a copy of that model. The question the case puts to it is the one
every model faces: whether a small network estimated on the joint loss earns a place among the econometric models under the same eligibility screens and the same ranking, and the answer is
reported whether it wins or not.

\subsubsection{The network}
For each day $t$ of the window, the features are built from returns before $t$ only. A volatility $\sigma_t$ is the RiskMetrics recursion of equation~\eqref{eq:mk-ewma} with $\lambda=0.94$ seeded on the variance of
the window's first 20 days, and a slow volatility $\sigma_t^{slow}$ is the same recursion with $\lambda=0.99$. The input vector has $2\cdot5+1=11$ entries: the five previous returns and their absolute values, each divided by
$\sigma_t$, and the log ratio of the fast to the slow volatility,
\begin{equation}
x_t=\Big(\tfrac{y_{t-1}}{\sigma_t},\dots,\tfrac{y_{t-5}}{\sigma_t},\ \tfrac{|y_{t-1}|}{\sigma_t},\dots,\tfrac{|y_{t-5}|}{\sigma_t},\ \ln\tfrac{\sigma_t}{\sigma_t^{slow}}\Big).
\label{eq:mk-nn-feat}
\end{equation}
The inputs are standardised by the mean and standard deviation of the training rows. A hidden layer of $H=\mkNnHidden{}$ units with the hyperbolic tangent feeds two outputs $o_1,o_2$, which two softplus functions turn into a VaR below
zero and an ES below the VaR,
\begin{equation}
h_t=\tanh\!\left(W_1\tilde x_t+b_1\right),\qquad o_t=W_2h_t+b_2,\qquad
v_t^{z}=-\mathrm{softplus}(o_{t,1}),\qquad e_t^{z}=v_t^{z}-\mathrm{softplus}(o_{t,2}),
\label{eq:mk-nn-out}
\end{equation}
where $\tilde x_t$ is the standardised input and $\mathrm{softplus}(x)=\ln(1+e^x)$. The construction guarantees $e_t^z\le v_t^z<0$ for every input, the domain of FZ0, without a penalty. The network has
\mkNnParams{} parameters: $11H+H$ in the first layer and $2H+2$ in the second. The target is the standardised return $z_t=y_t/\sigma_t$, and the loss is the mean FZ0 of the training rows,
\begin{equation}
\mathcal L(\theta)=\frac1{n_{tr}}\sum_{t\in\mathrm{train}}L_{FZ0}\big(z_t,\,v^z_t(\theta),\,e^z_t(\theta);\alpha\big),
\label{eq:mk-nn-loss}
\end{equation}
trained separately for each tail level. The forecast for the day after the window is the volatility times the network's standardised pair, $(\widehat v,\widehat e)=\sigma_{t}\,(v^z_t,e^z_t)$. The
following statement is what makes training on the standardised problem legitimate.

\begin{proposition}[Scale covariance of the FZ0 loss]
\label{prop:mk-scale}
For $c>0$, $L_{FZ0}(cY,cv,ce;\alpha)=L_{FZ0}(Y,v,e;\alpha)+\ln c$. Consequently, if $\sigma_t$ is known before day $t$ and $(v,e)$ minimises the expected FZ0 loss of the standardised return $Y/\sigma_t$ given $\mathcal F_{t-1}$, then $(\sigma_tv,\sigma_te)$ minimises the
expected FZ0 loss of $Y$ given $\mathcal F_{t-1}$.
\end{proposition}
\begin{proof}
In equation~\eqref{eq:mk-fz0} the indicator $\indic\{cY\le cv\}$ equals $\indic\{Y\le v\}$, and the ratios $(cv-cY)/(ce)$ and $cv/(ce)$ equal $(v-Y)/e$ and $v/e$, so that only the logarithm changes, $\ln(-ce)=\ln c+\ln(-e)$. Given
$\mathcal F_{t-1}$ the constant $\ln\sigma_t$ does not depend on $(v,e)$, so the minimiser of $\mathbb E\,L_{FZ0}(Y,\sigma_tv,\sigma_te)=\mathbb E\,L_{FZ0}(Y/\sigma_t,v,e)+\ln\sigma_t$ over $(v,e)$ is the minimiser for the
standardised variable, scaled by $\sigma_t$.
\end{proof}

Because the loss shifts by $\ln c$ under a change of scale, the average loss of the network on standardised returns, which the fit reports as its training loss, is not comparable with the out-of-sample
average FZ0 of any model on the original returns, and the product never puts the two side by side.

\subsubsection{Training, parameters and the choice of the hyperparameters}
The forecast of day $t$ trains on the estimation window $y_{t-1000},\dots,y_{t-1}$ only, and the day itself and everything after it enter no feature, no target and no standardisation constant. The first 60 days of the
window only feed the volatility recursions and the lags, so that the window of \mkNnWin{} days gives $\mkNnTrainRows{}$ training rows, and a window with fewer than 100 rows or no volatility is a failed fit. The weights of
the first layer are drawn from a normal law with standard deviation $1/\sqrt{11}$, those of the second layer from a normal law with standard deviation 0.1, and the output bias starts at the inverse softplus of the
normal quantile $|\Phi^{-1}(\alpha)|$ and of $0.15|\Phi^{-1}(\alpha)|$, so that the untrained network forecasts a Gaussian-like pair. Training is full-batch Adam \citep{market:kingma2015adam} for \mkNnEpochs{} epochs with a learning rate of \mkNnLr{} and a weight decay of
\mkNnWd{}, from \mkNnStarts{} seeded starts, and the start with the lowest final training loss is kept; the pair is refitted every \mkNnRefit{} forecast days and the weights are kept between refits, with the features
running over the current window. A fit that fails keeps the previous weights and is counted; none failed in the windows of C11 (\mkNnFailures{} failures in \mkNnRefits{} refits over the two levels). The softplus outputs make the loss
differentiable, and the gradient is computed by automatic differentiation of the loss only; the forward pass of a forecast is plain array arithmetic.

\Cref{alg:mk-net} states the procedure. The hidden units and the weight decay are chosen on the selection span of 1995 to 2003 like every other hyperparameter, one at a time from the base of 8 units and a decay of 0.001, by the lowest summed average FZ0 loss of the network itself
over the equity and the rates series: the units among 4, 8 and 16 and the weight decay among 0.001, 0.01 and 0.1, with the selected values \mkSelNnHiddenUnits{} units and a decay of \mkSelNnWeightDecay{}
(\cref{tab:market-selection}). The score rises with the number of hidden units and falls with the decay, so the grid is ordered toward the smallest, most regularised network; both
selected values lie at the edge of their grids, so a larger decay was not tried, and the selected pair was never scored jointly. The epochs, the learning rate, the number of starts and the refit interval are fixed by the compute limit and are not tuned.

\begin{algorithm}[htbp]
\caption{Fit and forecast of the quantile network for one forecast day and one tail level.}
\label{alg:mk-net}
\begin{algorithmic}[1]
\Require window $y_{t-N},\dots,y_{t-1}$ with $N=1000$, level $\alpha$, hyperparameters $(H,\text{epochs},\text{rate},\text{decay},\text{starts})$, seed
\State compute $\sigma$ and $\sigma^{slow}$ by the two recursions; build the rows $x_s$ for $s=60,\dots,t$; standardise by the training mean and standard deviation
\State targets $z_s\gets y_s/\sigma_s$ for the training rows $s<t$
\For{each start $j=1,\dots,\text{starts}$}
\State initialise $W_1,b_1,W_2,b_2$ from the seeded generator with the output bias at the inverse softplus of $|\Phi^{-1}(\alpha)|$ and $0.15|\Phi^{-1}(\alpha)|$
\For{epoch $=1,\dots,\text{epochs}$}
\State $v^z,e^z\gets$ forward pass of equation~\eqref{eq:mk-nn-out}; $\mathcal L\gets$ mean FZ0 of equation~\eqref{eq:mk-nn-loss}; stop the start if $\mathcal L$ is not finite
\State one Adam step on $(W_1,b_1,W_2,b_2)$ with the gradient of $\mathcal L$
\EndFor
\State record the final training loss of the start
\EndFor
\State keep the weights of the start with the lowest final loss
\State \Return $\sigma_t\,v^z_t$ and $\sigma_t\,e^z_t$ computed from the row of day $t$
\end{algorithmic}
\end{algorithm}

\subsubsection{Assumptions, complexity and implementation}
The network assumes that the standardised return depends on the recent past only through the five lags and the two volatilities; that a window of \mkNnTrainRows{} days carries enough information to fit a
tail quantile; and that the FZ0 loss is a good surrogate for the quality of the forecast pair. The second assumption is the weak one. At the 99 percent level a training window holds $\mkNnTailA{}$ tail observations
in expectation, and $\mkNnTailB{}$ at the 97.5 percent level; the objective at the lower level is determined by a handful of days, and the network can place its VaR path on them. One forward and backward pass costs
$O(n_{tr}\,H\,(11+2))$, so a refit costs $O(E\,S\,n_{tr}\,H)$ for $E$ epochs and $S$ starts, which for the parameters of the case is $E\,S\,n_{tr}=\mkNnCost{}$ row passes; a forecast costs $O(N)$ for the volatility
recursions and $O(H)$ for the forward pass. The refits are $T/\mkNnRefit{}$ per window and level.

The engine has no neural estimator and its FZ0 has no gradient, so the network is the pipeline's own code, in \mkfn{data-pipeline/pipeline/cases/c11/quantile_net.py}: \mkfn{features} builds the causal rows,
\mkfn{net_loss} is the FZ0 as an expression with a gradient, \mkfn{fit_net} trains the starts and \mkfn{predict} returns the pair of the day after the window. The gradient is computed with PyTorch 2.14.1 (a CPU wheel);
the loss is checked against the engine's \mkfn{fz0_loss} by a test of the case. The network goes through the same eligibility screens, the same FZ0 champion rule, the same Diebold--Mariano comparison and the same
model confidence set as the other estimators (\cref{sec:m19}). The acceptance criterion of the product is that the network is trained on each window's training span only, with no look-ahead, that its VaR and ES
pass the same eligibility tests and FZ0 ranking as the other models, and that its place is reported whether it wins or not. Its behaviour is read in \cref{sec:results-market}.

% ----------------------------------------------------------------------------------------------------------------------------------
\subsection{Results}
\label{sec:results-market}

\subsubsection{Case C11: market VaR and ES over six windows}
\Cref{tab:market-windows} gives the outcome of each window. In the control window of \mkStartCalm{} to \mkEndCalm{}, \mkEligCalm{} of the \mkNModels{} models are eligible and the champion is GJR-t at 99 percent and the EWMA at
97.5 percent; of the thirty-six Kupiec tests, \mkRedKupCalm{} are red. In the equity crisis of \mkStartGfc{} to \mkEndGfc{} no model is eligible: all thirty-six Kupiec tests, all thirty-six conditional coverage tests and
all thirty-six dynamic quantile tests are red, and the Acerbi--Szekely test is red for \mkRedAzGfc{} of the \mkEvAzGfc{} model and level pairs that have an ES. The models closest to coverage, among them the extreme-value and the filtered historical simulations,
have \mkExcEvtGfc{} exceedances of the 99 percent VaR against \mkExpGfc{} expected, and the delta-normal model has \mkExcDnGfc{}, about eleven times the expected count (\cref{tab:market-exceed}, \cref{fig:market-hitrate}). The window has
no champion, which is a result: when every candidate is rejected, a ranking under a consistent score still exists (\cref{fig:market-fz0}) but names no model that the coverage evidence supports. In the rates window,
\mkEligRates{} models are eligible and the EWMA is the champion at both levels; in the pandemic window \mkEligCovid{} are eligible over only \mkDaysCovid{} forecast days, with the filtered historical simulation the champion at 99 percent and the conditional
extreme-value model at 97.5; in the euro window \mkEligFx{} are eligible. In the last window, \mkEligMixed{} of \mkNModels{} models are eligible, but the expected number of exceedances at 99 percent is \mkExpMixed{} and the independence,
conditional coverage and duration tests could not be evaluated for most models (\cref{tab:market-reds}); the window has \mkDaysMixed{} days, and a pass says almost nothing.

\begin{table}[htbp]
\centering
\footnotesize
\setlength{\tabcolsep}{3.0pt}
\caption{Case C11, one row per window: the book, the number of forecast days, the number of the eighteen models that pass the coverage, conditional coverage and expected shortfall tests at both tail levels, the champion at each level with its average FZ0 loss, and the size of the 90 percent model confidence set at each level.}
\label{tab:market-windows}
\begin{tabular}{llrrp{3.2cm}p{3.2cm}rr}
\toprule
Window & Book & Days & Eligible & Champion at 1\% & Champion at 2.5\% & In set 1\% & In set 2.5\% \\
\midrule
Calm 2004--06 & mixed & 750 & 11 & GJR-t (-4.521) & EWMA (-4.659) & 13 & 13 \\
Equity 2007--09 & equity & 504 & 0 & none & none & 13 & 13 \\
FX 2008--09 & fx & 251 & 11 & ES-CAViaR (-3.557) & Cond. EVT (-3.716) & 14 & 14 \\
Equity 2020 & equity & 104 & 4 & Filtered HS (-1.986) & Cond. EVT (-2.364) & 15 & 14 \\
Rates 2022 & rates & 249 & 4 & EWMA (-3.811) & EWMA (-4.002) & 15 & 13 \\
Mixed 2023 & mixed & 64 & 18 & GJR-t (-4.439) & GJR-t (-4.623) & 2 & 4 \\
\bottomrule
\end{tabular}
\end{table}


\begin{table}[htbp]
\centering
\scriptsize
\setlength{\tabcolsep}{2.0pt}
\caption{Exceedances of the 99 percent value at risk of each model in each window of case C11, with the number expected under correct coverage in the first row; a dagger marks a model whose Kupiec test is red at the 5 percent level in that window. The daily forecasts are strictly out of sample.}
\label{tab:market-exceed}
\begin{tabular}{lrrrrrr}
\toprule
Model & Calm 2004--06 & Equity 2007--09 & FX 2008--09 & Equity 2020 & Rates 2022 & Mixed 2023 \\
\midrule
Expected & 7.5 & 5.0 & 2.5 & 1.0 & 2.5 & 0.6 \\
HS 250 & \relax 9 & \relax 19$^{\dagger}$ & \relax 10$^{\dagger}$ & \relax 8$^{\dagger}$ & \relax 8$^{\dagger}$ & \relax 0 \\
HS 500 & \relax 3 & \relax 27$^{\dagger}$ & \relax 13$^{\dagger}$ & \relax 10$^{\dagger}$ & \relax 8$^{\dagger}$ & \relax 0 \\
Age-weighted HS & \relax 10 & \relax 16$^{\dagger}$ & \relax 6 & \relax 5$^{\dagger}$ & \relax 7$^{\dagger}$ & \relax 0 \\
Delta-normal & \relax 0$^{\dagger}$ & \relax 57$^{\dagger}$ & \relax 23$^{\dagger}$ & \relax 16$^{\dagger}$ & \relax 15$^{\dagger}$ & \relax 0 \\
Cornish-Fisher & \relax 1$^{\dagger}$ & \relax 26$^{\dagger}$ & \relax 14$^{\dagger}$ & \relax 8$^{\dagger}$ & \relax 4 & \relax 0 \\
EWMA & \relax 10 & \relax 16$^{\dagger}$ & \relax 3 & \relax 7$^{\dagger}$ & \relax 6 & \relax 0 \\
GARCH-t & \relax 9 & \relax 16$^{\dagger}$ & \relax 3 & \relax 3 & \relax 6 & \relax 0 \\
GJR-t & \relax 9 & \relax 14$^{\dagger}$ & \relax 3 & \relax 3 & \relax 6 & \relax 0 \\
EGARCH-t & \relax 10 & \relax 20$^{\dagger}$ & \relax 3 & \relax 3 & \relax 6 & \relax 0 \\
Filtered HS & \relax 8 & \relax 13$^{\dagger}$ & \relax 3 & \relax 2 & \relax 5 & \relax 0 \\
Cond. EVT & \relax 8 & \relax 13$^{\dagger}$ & \relax 3 & \relax 2 & \relax 5 & \relax 0 \\
CAViaR-SAV & \relax 4 & \relax 13$^{\dagger}$ & \relax 5 & \relax 4$^{\dagger}$ & \relax 6 & \relax 0 \\
CAViaR-AS & \relax 5 & \relax 15$^{\dagger}$ & \relax 4 & \relax 4$^{\dagger}$ & \relax 6 & \relax 0 \\
ES-CAViaR & \relax 4 & \relax 13$^{\dagger}$ & \relax 5 & \relax 4$^{\dagger}$ & \relax 6 & \relax 0 \\
GAS-FZ 1F & \relax 4 & \relax 20$^{\dagger}$ & \relax 11$^{\dagger}$ & \relax 6$^{\dagger}$ & \relax 10$^{\dagger}$ & \relax 0 \\
GAS-FZ 2F & \relax 3 & \relax 26$^{\dagger}$ & \relax 19$^{\dagger}$ & \relax 7$^{\dagger}$ & \relax 8$^{\dagger}$ & \relax 0 \\
Hybrid GARCH-FZ & \relax 3 & \relax 27$^{\dagger}$ & \relax 16$^{\dagger}$ & \relax 6$^{\dagger}$ & \relax 9$^{\dagger}$ & \relax 0 \\
Quantile network & \relax 22$^{\dagger}$ & \relax 16$^{\dagger}$ & \relax 5 & \relax 4$^{\dagger}$ & \relax 11$^{\dagger}$ & \relax 1 \\
\bottomrule
\end{tabular}
\end{table}


\begin{figure}[htbp]
\centering
\includegraphics[width=0.95\textwidth]{market-hitrate.pdf}
\caption{Exceedances of the 99 percent value at risk divided by the number expected under correct coverage, for the seventeen ladder models (grey circles) and the quantile network (diamond), in each of the six windows of
case C11. The grey boxes are the central 95 percent range of the binomial law of the count for the length of the window, so a point above its box is a model with too many exceedances; the dashed line is a ratio of one.}
\label{fig:market-hitrate}
\end{figure}

\Cref{fig:market-hitrate} puts the exceedance counts on one scale. In the crisis window all \mkAboveGfc{} models lie above the acceptance box, whose upper end is \mkBandHiGfc{} exceedances; in the pandemic window \mkAboveCovid{} models do, with the box ending
at \mkBandHiCovid{}; in the rates window \mkAboveRates{}, in the euro window \mkAboveFx{}, and in the calm control window one, the quantile network. The box is wide in the short windows: in the last window it ends at \mkBandHiMixed{}
exceedances and no model leaves it, which is the lack of power that the text describes. The figure does not show whether the exceedances cluster, which the independence and duration tests read, nor the size of the loss beyond the VaR, which the
expected shortfall test reads. The Basel traffic light adds the view over time: for the 250-day historical simulation at 99 percent in the crisis window, of the \mkTlN{} rolling windows of 250 days none is green, \mkTlAmber{} are amber and \mkTlRed{} are red, the largest
rolling count is \mkTlMax{} exceptions, and \mkTlFinalRed{} of the \mkNModels{} models end the window in the red zone (\cref{tab:market-traffic} gives the zones). The windows shorter than 250 forecast days have no light.

\begin{table}[htbp]
\centering
\scriptsize
\setlength{\tabcolsep}{2.5pt}
\caption{Number of red lights of each back-test over the model and tail-level pairs that the test could evaluate (at most thirty-six pairs: eighteen models at two levels), per window of case C11. A test is not evaluated when its statistic is undefined, for example with no exceedance or no expected shortfall forecast.}
\label{tab:market-reds}
\begin{tabular}{lrrrrrr}
\toprule
Test & Calm 2004--06 & Equity 2007--09 & FX 2008--09 & Equity 2020 & Rates 2022 & Mixed 2023 \\
\midrule
\relax Kupiec proportion of failures & 10/36 & 36/36 & 13/36 & 25/36 & 21/36 & 0/36 \\
\relax Christoffersen independence & 0/35 & 4/36 & 3/36 & 0/36 & 0/36 & 0/4 \\
\relax Christoffersen conditional coverage & 7/35 & 36/36 & 13/36 & 21/36 & 18/36 & 0/4 \\
\relax Duration test & 2/34 & 7/36 & 1/36 & 2/36 & 4/36 & 0/0 \\
\relax Dynamic quantile test & 4/36 & 36/36 & 20/36 & 24/36 & 25/36 & 0/36 \\
\relax Acerbi-Szekely $Z_2$ & 0/30 & 16/30 & 9/30 & 14/30 & 4/30 & 0/30 \\
\bottomrule
\end{tabular}
\end{table}


\Cref{tab:market-reds} reads the whole battery. In the crisis window the dynamic quantile test, the Kupiec test and the conditional coverage test are red for every pair; the independence test is red for only \mkRedIndGfc{} and the duration test for \mkRedDurGfc{}, so the
evidence of the crisis is mainly of the level: the independence test, which is conditional on the frequency, rejects for few pairs. In the control window the dynamic quantile test
is red for \mkRedDqCalm{} pairs and the independence test for none. The Acerbi--Szekely test is not evaluated for six pairs in every window, the three models without an ES at two levels.

\begin{figure}[htbp]
\centering
\includegraphics[width=0.95\textwidth]{market-fz0.pdf}
\caption{Average FZ0 loss of each model minus the lowest average of its window, at the 99 percent level (left) and the 97.5 percent level (right), in each window of case C11; bold figures mark models that pass every eligibility
test, a darker cell is a worse forecast of the pair, and cells marked n/a are the models without an expected shortfall forecast. The scale is clipped at the 90th percentile of the differences.}
\label{fig:market-fz0}
\end{figure}

\Cref{fig:market-fz0} shows what the comparison under the consistent score adds to the coverage tests. The models that win on the score are not always the models that
pass the tests: the lowest average FZ0 of the crisis window at 99 percent belongs to \mkBestAnyAGfc{}, which fails the Kupiec and the conditional coverage tests. The delta-normal model is the clearest loser in the stress windows: it is worse than the best model by
\mkDnGapAGfc{} in the crisis window and by \mkDnGapACovid{} in the pandemic window at 99 percent, against \mkDnGapACalm{} in the calm window. The
model confidence set at 90 percent contains \mkMcsMinA{} to \mkMcsMaxA{} of the \mkMcsNACalm{} models with an ES at 99 percent in the five longer windows, so in every one of them the score removes at most two models and cannot separate the
rest from one another; in the short last window the set contains only \mkMcsAMixed{} models, because with no exceedance the FZ0 differences between models are nearly deterministic across days, which shrinks the variance of the differences and the set with it.

\begin{figure}[htbp]
\centering
\includegraphics[width=0.95\textwidth]{market-rates2022.pdf}
\caption{The rates window of 2022 in case C11. Upper panel: the realised loss of the 10-year bond position (grey points), the 99 percent value at risk of the EWMA, the GJR-GARCH-t model and the quantile network, and the exceedances of
the network (crosses). Lower panel: the cumulative number of exceedances of the same three models against the number expected under correct coverage (grey line).}
\label{fig:market-rates2022}
\end{figure}

\Cref{fig:market-rates2022} shows one window in time. The EWMA, the champion of the window, has \mkExcEwmaRates{} exceedances against \mkExpRates{} expected; the GJR model has \mkExcGjrRates{}, and the network has \mkNnExcRates{}, the largest
count of the three. The upper panel shows that the network's VaR path moves by larger steps from day to day than the other two. The figure shows when the exceedances occur and not why; the
reading of the network's path is in the next subsubsection.

\begin{figure}[htbp]
\centering
\includegraphics[width=0.95\textwidth]{market-z2null.pdf}
\caption{Simulated null distribution of the Acerbi--Szekely $Z_2$ statistic (bars: share of ten thousand simulated paths) for the age-weighted historical simulation in the equity crisis at 99 percent (left) and in the rates window
at 97.5 percent (right), with the realised statistic (solid line) and the critical values at the 5 percent and 0.01 percent levels (dashed and dotted).}
\label{fig:market-z2null}
\end{figure}

\Cref{fig:market-z2null} shows what the expected shortfall test does. In the crisis window the realised statistic of the age-weighted simulation lies far below the 0.01 percent critical value, so the ES forecast is wrong by an amount that no
simulated path reaches; in the rates window the realised statistic sits just below the 5 percent critical value, an amber reading: it is in the 5 percent tail of the null but far from the 0.01 percent tail. The figure shows two of
the thirty pairs of model and level that the test evaluates in each window, and it conditions on the tail-consistent null of equation~\eqref{eq:mk-null}.

\subsubsection{The quantile network in the ladder}
\Cref{tab:market-nn} sets the network beside the ladder in every window. It is eligible in \mkNnEligWindows{} of \mkNWindows{} windows (the euro window and the last window), the champion in \mkNnChampA{} of them at 99 percent and in \mkNnChampB{} at 97.5 percent, and a member of the 90
percent model confidence set in \mkNnSetA{} windows at 99 percent and in \mkNnSetB{} at 97.5 percent. Its average FZ0 at 99 percent ranks \mkNnRankCalm{} of \mkNRankCalm{} in the calm window, \mkNnRankGfc{} in the crisis, \mkNnRankFx{} in the euro window, \mkNnRankCovid{} in the pandemic window,
\mkNnRankRates{} in the rates window and \mkNnRankMixed{} in the last window, among the models that forecast an ES. In the pandemic window its 97.5 percent loss, \mkNnFzBCovid{}, is the lowest of all ladder models, but it fails the Kupiec test and is
therefore not the champion. The hypothesis that a small network trained on the joint loss earns the champion's place over the econometric models is not supported by these windows; its membership of the confidence set says only that the windows
are too short to separate it from the better models.

\begin{table}[htbp]
\centering
\scriptsize
\setlength{\tabcolsep}{2.0pt}
\caption{The quantile network trained on the FZ0 loss against the ladder in each window of case C11: eligibility, exceedances of its 99 percent value at risk with the expected number, the Kupiec p-value, its average FZ0 loss, the champion's average FZ0 loss at each level and whether the network belongs to the 90 percent model confidence set.}
\label{tab:market-nn}
\begin{tabular}{lcrrrrcrrc}
\toprule
Window & Eligible & Exc. 1\% (exp.) & Kupiec $p$ & FZ0 1\% & Best 1\% & Set 1\% & FZ0 2.5\% & Best 2.5\% & Set 2.5\% \\
\midrule
Calm 2004--06 & no & 22 (7.5) & $<0.001$ & -4.286 & -4.521 & yes & -4.588 & -4.659 & yes \\
Equity 2007--09 & no & 16 (5.0) & $<0.001$ & -2.371 & -- & yes & -2.856 & -- & yes \\
FX 2008--09 & yes & 5 (2.5) & 0.164 & -3.527 & -3.557 & yes & -3.685 & -3.716 & yes \\
Equity 2020 & no & 4 (1.0) & 0.026 & -1.669 & -1.986 & yes & -2.368 & -2.364 & yes \\
Rates 2022 & no & 11 (2.5) & $<0.001$ & -3.445 & -3.811 & yes & -3.909 & -4.002 & yes \\
Mixed 2023 & yes & 1 (0.6) & 0.676 & -4.134 & -4.439 & no & -4.374 & -4.623 & no \\
\bottomrule
\end{tabular}
\end{table}


The failure has a visible correlate. In the four windows in which the network fails an eligibility test, its Kupiec test is red at 99 percent, with \mkNnExcCalm{} exceedances against \mkExpCalm{} in the calm window, \mkNnExcGfc{} against \mkExpGfc{} in the crisis, \mkNnExcCovid{}
against \mkExpCovid{} in the pandemic window and \mkNnExcRates{} against \mkExpRates{} in the rates window; in the calm window it is the only model above the acceptance box of \cref{fig:market-hitrate}. \Cref{tab:market-rough} measures how much its VaR path moves: the mean absolute relative daily change of the network's VaR is between
\mkRoughMinNn{} and \mkRoughMaxNn{} across the windows, against \mkRoughMinGjr{} to \mkRoughMaxGjr{} for the GJR model, \mkRoughMinEwma{} to \mkRoughMaxEwma{} for the EWMA and \mkRoughMinHs{} to \mkRoughMaxHs{} for the 250-day historical simulation. A VaR path that changes by
\mkRoughMinNn{} to \mkRoughMaxNn{} of its value from one day to the next, on average, against at most \mkRoughMaxGjrOut{} for the GJR model outside the pandemic window (\mkGjrRoughACovid{} inside it), is not a quantile that tracks a slowly moving volatility; it responds to the five lags it receives. The training window holds only $\mkNnTailA{}$ tail observations at 99 percent in expectation, the objective at that level is dominated by a few days, and the fitted
function reproduces them without carrying them to the next day. The results are consistent with this account without proving it: the network has \mkNnParams{} parameters on \mkNnTrainRows{} rows, and the regularisation grid was ordered toward the smallest network. A
test of the mechanism, for example the same network with more regularisation or on a longer window, is not in the committed runs.

\begin{table}[htbp]
\centering
\footnotesize
\setlength{\tabcolsep}{4.0pt}
\caption{Mean absolute relative change from one day to the next of the 99 percent value at risk forecast by the quantile network and by three reference models, in each window of case C11; a larger value is a more volatile forecast path.}
\label{tab:market-rough}
\begin{tabular}{lrrrr}
\toprule
Window & Quantile network & GJR-t & EWMA & HS 250 \\
\midrule
Calm 2004--06 & 0.161 & 0.026 & 0.015 & 0.002 \\
Equity 2007--09 & 0.220 & 0.058 & 0.018 & 0.003 \\
FX 2008--09 & 0.211 & 0.020 & 0.018 & 0.002 \\
Equity 2020 & 0.373 & 0.171 & 0.030 & 0.012 \\
Rates 2022 & 0.251 & 0.066 & 0.016 & 0.003 \\
Mixed 2023 & 0.241 & 0.058 & 0.011 & 0.002 \\
\bottomrule
\end{tabular}
\end{table}


\subsubsection{Size and power of the tests}
\Cref{tab:market-sizepower} and \cref{fig:market-sizepower} are the measured behaviour of the tests on data with a known truth. Under the null, the rejection rate of every test lies between \mkSizeMin{} and \mkSizeMax{} percent across the \mkNStudies{}
studies; the Christoffersen independence and conditional coverage tests over-reject at 1,500 days (\mkSizeOverInd{} and \mkSizeOverCc{} percent, Wilson intervals \mkSizeOverIndLo{} to \mkSizeOverIndHi{} and \mkSizeOverCcLo{} to \mkSizeOverCcHi{}), the dynamic
quantile test at 250 days and 2.5 percent sits at the lower edge of its interval above 5 percent, and the others hold the nominal 5 percent within their intervals. The tests are not equally powerful. Against the stated alternatives, the Acerbi--Szekely test rejects in \mkPowAz{} percent of samples and the Christoffersen independence test in \mkPowInd{} percent, the Kupiec test in \mkPowKup{} percent against a hit rate of 6.5 percent in place of 5, and the dynamic quantile test
with 250 days at 99 percent in only \mkPowDqShort{} percent of samples against a clustered alternative, rising to \mkPowDqLong{} percent with 750 days. The first figure is the lack of power of the window length that the regulation uses, and it explains why a model can pass the dynamic quantile test in a window of that length with no clear meaning. The
alternatives are single points; the power elsewhere is not measured.

\begin{figure}[htbp]
\centering
\includegraphics[width=0.95\textwidth]{market-sizepower.pdf}
\caption{Size under the null (left, with the 5 percent level as a dashed line) and power against one alternative (right) of the nine back-tests in the studies of case C22, with the Wilson 95 percent interval of each rate; the label gives the
test and the sample length and level of the study.}
\label{fig:market-sizepower}
\end{figure}

\subsubsection{Case C13: what the desk tests catch}
\Cref{tab:market-pla,tab:market-desk} and \cref{fig:market-plaplane,fig:market-detection} give the outcome of the 1{,}000 replications of each variant. With nothing planted the attribution test is amber or red in \mkAmbNone{} percent of replications (Wilson interval \mkAmbLoNone{}
to \mkAmbHiNone{} percent), with a mean Spearman correlation of \mkSpNone{} (standard deviation \mkSpSdNone{}) and a mean Kolmogorov--Smirnov distance of \mkKsNone{}, and the desk is never over its limit: the tests have no false alarms on this generator. A missing basis is caught in \mkAmbMissingBasis{} percent of
replications (interval \mkAmbLoMissingBasis{} to \mkAmbHiMissingBasis{}), red in \mkRedMissingBasis{} percent; the Spearman statistic flags \mkSpFlagMissingBasis{} percent and the Kolmogorov--Smirnov statistic \mkKsFlagMissingBasis{} percent, and the desk back-test \mkDeskFlagMissingBasis{} percent. An omitted factor is caught in \mkAmbOmittedFactor{} percent
(red in only \mkRedOmittedFactor{} percent), mainly by the Spearman statistic (\mkSpFlagOmittedFactor{} percent) and less by the Kolmogorov--Smirnov statistic (\mkKsFlagOmittedFactor{} percent), while the desk back-test catches \mkDeskFlagOmittedFactor{} percent. A P\&L dated one day off is red in \mkRedTimeMisalignment{} percent, through the Spearman statistic alone, and neither the Kolmogorov--Smirnov statistic nor the desk back-test sees it
because it changes neither the marginal distribution of the risk-theoretical P\&L nor the desk VaR.

\begin{table}[htbp]
\centering
\scriptsize
\setlength{\tabcolsep}{2.0pt}
\caption{Case C13, the P\&L attribution test on 1,000 replications of each variant: the mean and standard deviation of the Spearman correlation and of the Kolmogorov-Smirnov distance, the number of replications in each joint zone, and the share that is amber or red with its Wilson 95 percent interval in brackets.}
\label{tab:market-pla}
\begin{tabular}{lrrrrrr}
\toprule
Variant & Spearman, mean (sd) & KS, mean (sd) & Green & Amber & Red & Amber or red (\%) \\
\midrule
None (size) & 0.992 (0.004) & 0.033 (0.006) & 1000 & 0 & 0 & 0.0 [0.0, 0.4] \\
Missing basis & 0.705 (0.073) & 0.120 (0.028) & 40 & 385 & 575 & 96.0 [94.6, 97.0] \\
Stale volatility & 0.979 (0.005) & 0.041 (0.008) & 1000 & 0 & 0 & 0.0 [0.0, 0.4] \\
Wrong correlation & 0.986 (0.004) & 0.141 (0.014) & 0 & 87 & 913 & 100.0 [99.6, 100.0] \\
Omitted factor & 0.772 (0.031) & 0.089 (0.016) & 99 & 858 & 43 & 90.1 [88.1, 91.8] \\
Time misalignment & -0.004 (0.069) & 0.033 (0.006) & 0 & 0 & 1000 & 100.0 [99.6, 100.0] \\
\bottomrule
\end{tabular}
\end{table}


\begin{figure}[htbp]
\centering
\includegraphics[width=0.95\textwidth]{market-plaplane.pdf}
\caption{The six variants of case C13 in the plane of the two attribution statistics: the mean Spearman correlation and the mean Kolmogorov--Smirnov distance over 1,000 replications, with one standard deviation as bars. The grey box is the green
zone of the framework, and the dashed lines are the red boundaries of the Spearman correlation and of the distance.}
\label{fig:market-plaplane}
\end{figure}

Two variants escape the attribution test. A model that overstates the correlation by 0.3 has a mean Spearman correlation of \mkSpWrongCorrelation{}, never flagged (\mkSpFlagWrongCorrelation{} percent), but a mean distance of \mkKsWrongCorrelation{}, flagged in \mkKsFlagWrongCorrelation{} percent of replications, and the desk
back-test sees it too: the desk is over its 99 percent limit in \mkOverNNWrongCorrelation{} percent of replications and over its 97.5 percent limit in \mkOverNFWrongCorrelation{}. A model that scales the shocks with the volatility of 20 days before is flagged by neither the attribution test (\mkAmbStaleVolatility{} percent amber or red, interval \mkAmbLoStaleVolatility{}
to \mkAmbHiStaleVolatility{}) nor, in practice, the desk back-test (\mkOverNNStaleVolatility{} percent over the 99 percent limit). The closed form of equation~\eqref{eq:mk-pearson} explains it: the population Pearson correlation of the two series is still \mkPearsonStaleVolatility{}, the
measured mean Spearman correlation is \mkSpStaleVolatility{}, and the marginal distribution of the lagged series is that of the unlagged one, so that a statistic that compares marginal distributions has nothing to detect. The deficiency is real, because the desk VaR lags the true scale, and the
attribution test does not see it. The figure shows means and not the spread of the replications, which \cref{tab:market-pla} gives, and it holds for this desk of one stylised book, not for any desk of a bank.

\begin{table}[htbp]
\centering
\scriptsize
\setlength{\tabcolsep}{2.0pt}
\caption{Case C13, the desk back-test on 1,000 replications of each variant: the exceptions of the hypothetical P\&L against the risk model's value at risk at the 99 and 97.5 percent levels (mean and standard deviation), the percentage of replications over the limit of the framework at each level, and the number of replications flagged by the Spearman statistic, by the Kolmogorov-Smirnov statistic and by the desk back-test at either level.}
\label{tab:market-desk}
\begin{tabular}{lrrrrrrr}
\toprule
Variant & Exc. 99\% & Over 99\% & Exc. 97.5\% & Over 97.5\% & Spearman & KS & Desk \\
\midrule
None (size) & 2.6 (1.6) & 0.0 & 6.5 (2.5) & 0.0 & 0 & 0 & 0 \\
Missing basis & 14.2 (5.8) & 56.5 & 22.2 (7.1) & 13.0 & 919 & 867 & 565 \\
Stale volatility & 5.0 (2.4) & 0.5 & 9.2 (3.3) & 0.0 & 0 & 0 & 5 \\
Wrong correlation & 19.3 (4.2) & 95.3 & 28.8 (4.9) & 36.7 & 0 & 1000 & 954 \\
Omitted factor & 8.0 (2.8) & 6.4 & 14.9 (3.8) & 0.0 & 805 & 448 & 64 \\
Time misalignment & 2.7 (1.6) & 0.0 & 6.5 (2.6) & 0.0 & 1000 & 0 & 0 \\
\bottomrule
\end{tabular}
\end{table}


\begin{figure}[htbp]
\centering
\includegraphics[width=0.95\textwidth]{market-detection.pdf}
\caption{Share of the 1,000 replications of each variant of case C13 flagged by each test, with the Wilson 95 percent interval: the Spearman statistic (amber or red), the Kolmogorov--Smirnov statistic (amber or red) and the desk back-test (over the limit at either level). The labels give
the percentage.}
\label{fig:market-detection}
\end{figure}

\Cref{fig:market-detection} puts the three statistics side by side. No statistic dominates: the Spearman statistic is the only one that catches a one-day misalignment, the distance is the only one of the two attribution statistics that catches an overstated correlation, and the desk back-test catches the
overstated correlation and part of the missing basis (\mkDeskFlagMissingBasis{} percent) but almost none of the omitted factor (\mkDeskFlagOmittedFactor{} percent) or of the stale volatility (\mkDeskFlagStaleVolatility{} percent). The joint attribution zone flags four of the five deficiencies in at least \mkAmbOmittedFactor{} percent of
replications, the lowest of the four being the omitted factor, and the fifth, the stale volatility, in none of them; the desk back-test adds almost nothing for that one. The reading is conditional on the strengths that the generator plants, which are choices of the case and not estimates from any bank.
